\documentclass[../main.tex]{subfiles}
\begin{document}

% Nguồn số liệu của chương: repo SUMO (CLAUDE.md, pyproject.toml, docs/PROGRESS.md,
% docs/research/phase3_box_capacity.md, experiments/phase3/results_r21.csv,
% experiments/operator/task4b/b0_table.csv, src/thesis_sumo/**; cập nhật 2026-10-28: docs/research/phase3*.md,
% report/2026-10-27/data/campaign_specs.csv, experiments/campaigns/{P1,P2}/loop/iter_*.json,
% experiments/llm_logs/{P1,P2}/validity_log.csv, experiments/campaigns/REF_run.log; cập nhật 2026-10-31: report/2026-10-28, report/2026-10-29,
% experiments/campaigns/{REF_n29s5,AB3_*,AB_*}/runs/**/manifest.json, REF_n29s5/{B0,B1_manual}.pkl); cập nhật 2026-11-09: report/2026-11-08/README.md
% (REF_n1104s5, AB_cut, B2/P1/P2_n1108, AB_d1, AB_d1v). Các mục [TODO] cần bổ sung trước khi nộp.

Chương~\ref{chapter:Methodology} đã trình bày phương pháp ở mức thiết kế. Chương này mô tả phần mềm hiện thực phương pháp đó và các kết quả thực nghiệm đã có, để kiểm chứng rằng từng thành phần hoạt động đúng và để ghi lại những gì quy trình cho thấy trên dữ liệu đại diện. Chương gồm năm phần: kiến trúc phần mềm (Mục~\ref{section:architecture}), công cụ và kết quả triển khai (Mục~\ref{section:tools}), kiểm thử (Mục~\ref{section:testing}), thực nghiệm (Mục~\ref{section:experiments}) và cách tái lập (Mục~\ref{section:deploy}). Tại ngày 08/11/2026, toàn bộ mã của quy trình đã chạy được, từ mạng mô phỏng, bộ điều khiển TraCI, chỉ số, khung đánh giá đến vòng lặp LLM và phương án Optuna. Ba chiến dịch B2, P1 và P2 (P1, P2 với \texttt{qwen3.5:9b}) đã chạy trên mạng 04/11, và không đặc tả nào đủ điều kiện (Mục~\ref{subsection:campaigns_n1108}). Phương án H và đánh giá cuối chưa có kết quả (Bảng~\ref{table:3.1}).

Mọi kết quả định lượng trong chương được tính trên mạng và dữ liệu đại diện (\texttt{LITERATURE\_PROXY}) dựng từ tài liệu, không phải số đo tại nút. Giá trị đầu vào tạm thời in \tamthoi{màu đỏ} theo quy ước ở Mục~\ref{section:3.2}.

\section{Kiến trúc phần mềm}
\label{section:architecture}

\subsection{Đường ống và hai giai đoạn}

Hệ thống là một đường ống (pipeline) gồm các bước tách rời; mỗi bước đọc tệp đầu vào và ghi tệp đầu ra dạng văn bản (YAML, JSON, XML, CSV). Cách tổ chức này đáp ứng ba yêu cầu. Yêu cầu thứ nhất là \emph{tái lập}: mạng sinh lại từ dữ liệu gốc bằng script và tệp vá, và mọi lần chạy dùng seed cố định và ghi \texttt{manifest.json}. Yêu cầu thứ hai là \emph{thay được nguồn đặc tả}: B0, B1, Optuna và LLM đều chỉ tạo một đặc tả JSON cùng lược đồ, phần còn lại của đường ống không phân biệt nguồn. Yêu cầu thứ ba là \emph{dùng chung hàm chỉ số}: đầu ra SUMO được chuyển về lược đồ của dữ liệu quan sát, nên một bộ hàm tính được cho cả hai.

Thời gian chạy chia hai giai đoạn. Giai đoạn \emph{ngoại tuyến} gồm sinh hoặc soạn đặc tả, kiểm tra hợp lệ và biên dịch; LLM chỉ tham gia ở đây và chỉ tạo tệp JSON. Giai đoạn \emph{trực tuyến} là một lần chạy SUMO: SUMO nạp tham số tĩnh từ \texttt{vtypes.add.xml} khi khởi động, còn bộ điều khiển TraCI đọc \texttt{traci\_rules.json} và đánh giá luật ở mỗi bước. Không có lời gọi LLM nào trong lúc mô phỏng.

Một trình chạy (runner) duy nhất điều phối mỗi lần chạy. Nó xác định đầu vào theo quy ước tên (mạng theo tệp cấu hình hiện trường, tệp tuyến theo ngày, đặc tả đã biên dịch theo mã đặc tả), khởi động SUMO qua \texttt{libsumo} (chạy hàng loạt) hoặc \texttt{traci} (khi cần \texttt{sumo-gui}), gắn bộ điều khiển nếu đặc tả có luật, ghi \texttt{manifest.json} và chuyển đầu ra sang lược đồ quan sát. Các seed chạy song song, mỗi tiến trình một phiên SUMO.

\subsection{Luồng dữ liệu}

Bảng~\ref{table:4.2} mô tả chín bước của luồng dữ liệu. Bước 1--3 chạy một lần cho mỗi cấu hình hiện trường hoặc mỗi ngày dữ liệu. Bước 4--8 lặp lại cho mỗi đặc tả. Bước 9 đóng vòng lặp bằng cách đưa kết quả về bộ sinh đặc tả.

\begin{table}[H]
\centering
\small
\caption{Luồng dữ liệu giữa các gói.} \label{table:4.2}
\begin{tabular}{|p{0.6cm}|p{3.4cm}|p{3.6cm}|p{4.9cm}|}
\hline
\textbf{\#} & \textbf{Bước (gói)} & \textbf{Đầu vào} & \textbf{Đầu ra} \\ \hline
1 & Dựng mạng (\texttt{net}) & Bản xuất OSM, tệp cấu hình hiện trường & Mạng cơ sở, mạng hiện trường, tệp vá, bảng cạnh, sơ đồ mạng \\ \hline
2 & Dữ liệu mục tiêu (\texttt{proxy}) & Bảng giá trị mục tiêu, cấu hình hiện trường & Các bảng CSV theo lược đồ quan sát của từng ngày \\ \hline
3 & Nhu cầu (\texttt{demand}) & Mạng, số đếm rẽ 15 phút & Tệp tuyến của từng ngày \\ \hline
4 & Kiểm tra hợp lệ (\texttt{rules}) & Đặc tả JSON, miền tham số, cấu hình hiện trường, mạng & Danh sách lỗi có mã \\ \hline
5 & Biên dịch (\texttt{rules}) & Đặc tả hợp lệ & \texttt{vtypes.add.xml}, \texttt{traci\_\allowbreak rules.json} (nếu có luật) \\ \hline
6 & Chạy mô phỏng (\texttt{sim}, \texttt{runtime}) & Mạng, tệp tuyến, đặc tả đã biên dịch, cấu hình thí nghiệm, seed & Đầu ra SUMO, \texttt{sumo.log}, \texttt{interventions.csv}, \texttt{manifest.json} \\ \hline
7 & Chuyển lược đồ (\texttt{sim}) & Đầu ra SUMO & Các bảng CSV theo lược đồ quan sát \\ \hline
8 & Chỉ số và điểm (\texttt{metrics}, \texttt{evaluate}) & Bảng CSV mô phỏng và quan sát của các ngày hiệu chỉnh & M1--M10, G0, $R$, $S$, tính đủ điều kiện \\ \hline
9 & Sinh đặc tả kế tiếp (\texttt{llm}, \texttt{baselines}) & Kết quả đánh giá, ngân sách còn lại & Phản hồi cho LLM hoặc giá trị mục tiêu cho Optuna; đặc tả mới \\ \hline
\end{tabular}
\end{table}

Phụ thuộc giữa các gói đi một chiều. \texttt{config} và \texttt{schema\_observed} là nền, mọi gói khác dùng. \texttt{rules} và \texttt{runtime} dùng chung một mô-đun xác định kết nối đi sai làn, để bộ kiểm tra và bộ điều khiển hiểu giống nhau về các kết nối dùng được. \texttt{sim} gọi \texttt{runtime} khi đặc tả có luật TraCI. \texttt{proxy} sinh dữ liệu độc lập với đầu ra SUMO, nhưng dùng lại định nghĩa hình học và lược đồ của \texttt{sim} để hai phía cùng định nghĩa bảng. \texttt{metrics} chỉ đọc các bảng CSV theo lược đồ quan sát nên không phụ thuộc SUMO.

\subsection{Các gói và quyết định thiết kế}
\label{subsection:package_design}

Mã nguồn là một gói Python \texttt{thesis\_sumo}, chia thành các gói con theo trách nhiệm. Bảng~\ref{table:4.1} nêu trách nhiệm và quyết định thiết kế đáng chú ý của từng gói; mô-đun \texttt{evaluate} chứa khung đánh giá dùng chung cho mọi phương án. Chức năng của bộ kiểm tra, bộ biên dịch và bộ điều khiển đã mô tả ở Mục~\ref{section:3.4} và~\ref{section:3.5}.

\begingroup\small
\begin{longtable}{|p{2.3cm}|p{4.2cm}|p{7.5cm}|}
\caption{Các gói con của \texttt{thesis\_sumo} và quyết định thiết kế.}\label{table:4.1}\\
\hline
\textbf{Gói} & \textbf{Trách nhiệm} & \textbf{Quyết định thiết kế} \\ \hline
\endfirsthead
\hline
\textbf{Gói} & \textbf{Trách nhiệm} & \textbf{Quyết định thiết kế} \\ \hline
\endhead
\texttt{config}, \texttt{schema\_\allowbreak observed} & Đọc cấu hình hiện trường, thí nghiệm, LLM; định nghĩa bảng và cột của lược đồ quan sát & Trường bắt buộc rỗng gây lỗi \texttt{MissingFieldDataError} nêu tên trường, không điền mặc định. Tệp đo thực mang nhãn \texttt{MEASURED} (chưa có giá trị); tệp đại diện mang nhãn \texttt{LITERATURE\_PROXY}, mỗi giá trị kèm nguồn và mức bằng chứng. Chương trình không bao giờ tự chuyển từ tệp đo thực sang tệp đại diện \\ \hline
\texttt{net} & Sinh mạng từ OSM, áp tệp vá, xuất bảng cạnh và sơ đồ mạng & Mạng cơ sở qua hai lượt \texttt{netconvert}: lượt đầu cắt theo bán kính và giữ cạnh cho ô tô, xe máy; lượt sau giữ bốn nhánh và đơn giản hóa hình học. Mạng hiện trường chạy lại \texttt{netconvert} tới khi điểm tách cạnh đúng khoảng cách (tối đa 5 lượt), rồi thêm lượt cuối đặt lại điểm chờ của xe rẽ trái trái phép. Việc tách cạnh làm vạch dừng dịch từ $-0{,}70$ đến $+1{,}61$~m tùy nhánh; độ dịch ghi vào một tệp để các bảng quan sát dùng đúng vị trí. Mạng 04/11 có 28 liên kết đèn: 24 liên kết lòng đường (0--23) và 4 liên kết vỉa hè (24--27); 10 liên kết chỉ cho xe vi phạm từ lane 0, 1 của mạng 29/10 (28--37) đã bỏ. Một lệnh dựng toàn bộ \\ \hline
\texttt{demand} & Sinh tuyến ứng viên và tệp tuyến từ số đếm rẽ 15 phút & Dùng công cụ \texttt{routeSampler} của SUMO \\ \hline
\texttt{proxy} & Sinh năm ngày dữ liệu đại diện PROXY\_D1--PROXY\_D5 từ bảng giá trị mục tiêu & Mỗi ngày gồm mười bảng CSV, một manifest và tệp tuyến. Không đọc đầu ra SUMO. Ba bảng nhóm F2 dùng luồng số ngẫu nhiên riêng để các bảng cũ không đổi \\ \hline
\texttt{rules} & Lược đồ đặc tả, bộ kiểm tra hợp lệ, bộ biên dịch & Bộ kiểm tra trả cả danh sách lỗi (tám mã), nội dung nêu rõ khoảng hoặc tên được phép để đưa nguyên văn vào phản hồi. Bộ biên dịch gọi bộ kiểm tra trước; đặc tả sai gây \texttt{InvalidSpecError} và không ghi tệp. Mọi tham số, kể cả bằng mặc định, được ghi tường minh. Tệp sinh ra khai báo XSD của SUMO, vì thử nghiệm cho thấy nếu không khai báo, thuộc tính sai tên bị bỏ qua không cảnh báo. Nguyên mẫu \texttt{custom1} (\texttt{custom2}) nhận trước giá trị mặc định của xe máy (ô tô), vì vClass tùy biến không có mặc định này. Mỗi luật mang danh sách mã vType thành viên, vì xe lấy từ một phân phối báo mã của vType thành viên chứ không phải của phân phối \\ \hline
\texttt{runtime} & Bộ điều khiển TraCI và sáu luật & Mỗi luật là lớp con của \texttt{Rule}, đăng ký trong \texttt{RULE\_REGISTRY}. Mỗi luật có bộ sinh số ngẫu nhiên riêng, khởi tạo từ seed, chỉ số và loại luật; xe duyệt theo thứ tự mã, nên \texttt{libsumo} và \texttt{traci} cho cùng chuỗi rút và cùng nhật ký. Mọi lệnh đi qua lớp bao \texttt{Commands}: mỗi xe một can thiệp, chỉ xóa đúng bit được nêu của speed mode, kiểm tra lại giá trị SUMO báo về. Một khoảng can thiệp chỉ đóng khi xe đã đi qua nút hoặc rời vùng xét, để va chạm trong nút của xe vượt đèn đỏ vẫn tính là liên quan can thiệp \\ \hline
\texttt{sim} & Ghi cấu hình SUMO, chạy, ghi manifest, chuyển đầu ra & \texttt{run.sumocfg} do chính SUMO ghi (tùy chọn lưu cấu hình), nên tùy chọn sai tên bị từ chối. Luôn bật \texttt{collision.check-junctions}. Ba trong mười bảng (chuỗi thời gian hàng chờ, sự kiện tìm khoảng trống, độ trễ xe rẽ phải) chỉ suy từ quỹ đạo, detector và tuyến, không đọc nhật ký can thiệp, để cùng định nghĩa áp dụng được cho dữ liệu video \\ \hline
\texttt{metrics} & Chỉ số M1--M10, G0, hàm khoảng cách, điểm tổng & Mỗi chỉ số một mô-đun, cộng một mô-đun dùng chung (gán chu kỳ đèn theo mọi mốc bắt đầu xanh, xác định xe rẽ phải liên tục và các hướng có mũi tên riêng để loại khỏi chỉ số vi phạm). Hàm \texttt{compare} trả $S$, $S_{\text{sourced}}$, $S_{\text{assumed}}$, khoảng cách và $R$ theo khóa, cùng các dòng tóm tắt không chứa tên ngày \\ \hline
\texttt{evaluate} & Ngân sách theo chiến dịch, đánh giá một đặc tả, chọn đặc tả tốt nhất & \texttt{RunBudget} giữ ngân sách trong \texttt{budget.json}, đọc lại từ đĩa, trừ dưới khóa tệp. Đặc tả không hợp lệ không tốn lần chạy; mã trùng trong một chiến dịch bị từ chối; ngày đánh giá bị chặn trừ khi gọi tường minh từ bước đánh giá cuối \\ \hline
\texttt{llm} & Gọi LLM, prompt \texttt{v001}--\texttt{v005}, phản hồi, vòng lặp P1/P2 & Hai backend: \texttt{openai} (yêu cầu khai báo nhiệt độ và mức lượng tử hóa) và \texttt{file} (ghi yêu cầu ra tệp và chờ tệp phản hồi, để chạy vòng lặp khi máy chủ LLM nằm ở máy khác). Phát lại nhật ký có kiểm định danh LLM, để câu trả lời của một mô hình không bị ghi nhận là của mô hình khác. \texttt{run\_loop} lưu trạng thái chiến dịch và từ chối chạy tiếp nếu trạng thái đổi \\ \hline
\texttt{baselines} & Optuna cho B2 và H & Nghiên cứu Optuna lưu trong tệp nhật ký, chạy tiếp không trừ lại ngân sách, từ chối nếu ngày, seed hoặc cấu trúc khác. Đặc tả tốt nhất chọn bằng cùng \texttt{select\_best} với vòng LLM, mang \texttt{provenance.model} là \texttt{optuna-tpe} \\ \hline
\end{longtable}
\endgroup

\subsection{Thiết kế lớp của bộ điều khiển}

Lớp cơ sở \texttt{Rule} giữ tham số luật, tập mã vType áp dụng, thông tin nút, lớp bao lệnh, bộ sinh số ngẫu nhiên và một từ điển các can thiệp đang diễn ra. Mỗi phần tử của từ điển ghi mã xe, hành động, thời điểm bắt đầu, speed mode gốc (nếu đã đổi) và cờ cho biết tốc độ đang bị ghi đè. Phương thức trừu tượng \texttt{on\_step} nhận ảnh chụp trạng thái của bước và trả về danh sách can thiệp đã có kết quả trong bước đó. Lớp cơ sở cung cấp năm thao tác dùng chung: kiểm tra xe thuộc phạm vi luật và đang rảnh; bắt đầu can thiệp; hoàn tác ghi đè (trả tốc độ cho SUMO, khôi phục speed mode); kết thúc can thiệp kèm kết quả và lý do; đóng mọi can thiệp còn dang dở khi lần chạy kết thúc, với lý do \texttt{run\_ended}. Khi SUMO báo speed mode khác giá trị vừa đặt, lớp bao lệnh báo \texttt{ModeNotRestoredError}.

Kết quả của một can thiệp là một bản ghi bất biến gồm bảy trường: (i) thời điểm phát lệnh; (ii) mã xe; (iii) loại luật; (iv) hành động; (v) cờ thực thi (phản ứng quan sát được của SUMO, định nghĩa riêng cho từng luật); (vi) lý do không thực thi (rỗng nếu đã thực thi); và (vii) thời điểm xác định kết quả.
Mỗi bản ghi là một dòng của \texttt{interventions.csv}.

\section{Công cụ và kết quả triển khai}
\label{section:tools}

\subsection{Thư viện và công cụ}

Bảng~\ref{table:4.3} liệt kê công cụ và thư viện. Cột ``Yêu cầu'' là ràng buộc tối thiểu khai báo trong cấu hình gói; cột ``Đã cài'' là phiên bản trong môi trường ảo của máy chính, ghi ngày 28/10/2026.

\begin{table}[H]
\centering
\small
\caption{Thư viện và công cụ sử dụng.} \label{table:4.3}
\begin{tabular}{|p{4.6cm}|p{4.4cm}|p{2.0cm}|p{2.0cm}|}
\hline
\textbf{Mục đích} & \textbf{Công cụ} & \textbf{Yêu cầu} & \textbf{Đã cài} \\ \hline
Mô phỏng vi mô, dựng mạng, sinh tuyến & SUMO (\texttt{sumo}, \texttt{sumo-gui}, \texttt{netconvert}, \texttt{routeSampler})~\cite{lopez2018sumo} & 1.27.0 & 1.27.0 \\ \hline
Điều khiển SUMO từ Python, chạy hàng loạt & \texttt{libsumo}, \texttt{sumo-data} & 1.27.0 & 1.27.0 \\ \hline
Điều khiển SUMO qua giao diện đồ họa, đọc mạng & \texttt{traci}, \texttt{sumolib} (lấy từ thư mục cài đặt SUMO) & như SUMO & 1.27.0 \\ \hline
Ngôn ngữ lập trình & Python (môi trường ảo) & $\geq$ 3.10 & 3.12.10 \\ \hline
Đọc cấu hình & \texttt{pyyaml} & $\geq$ 6.0 & 6.0.3 \\ \hline
Tính toán số & \texttt{numpy} & $\geq$ 1.26 & 2.5.4 \\ \hline
Xử lý bảng & \texttt{pandas} & $\geq$ 2.1 & 3.1.0 \\ \hline
Thống kê & \texttt{scipy} & $\geq$ 1.11 & 1.18.1 \\ \hline
Đọc và ghi XML & \texttt{lxml} & $\geq$ 5.0 & 6.1.3 \\ \hline
Vẽ hình & \texttt{matplotlib} & $\geq$ 3.8 & 3.11.2 \\ \hline
Kiểm tra đặc tả theo lược đồ JSON & \texttt{jsonschema} & $\geq$ 4.20 & 4.26.0 \\ \hline
Chuyển tọa độ kinh vĩ độ sang tọa độ mạng & \texttt{pyproj} & $\geq$ 3.6 & 3.8.0 \\ \hline
Kiểm thử tự động & \texttt{pytest} & $\geq$ 8.0 & 9.1.1 \\ \hline
Gọi LLM qua giao diện tương thích OpenAI & \texttt{openai} & $\geq$ 1.30 & 3.40.0 \\ \hline
Tối ưu tham số B2, H~\cite{akiba2019optuna} & \texttt{optuna} & $\geq$ 3.5 & 5.0.0 \\ \hline
Chạy LLM cục bộ & Ollama; mô hình \texttt{qwen3.5:9b}, Q4\_K\_M & -- & Ollama 0.40.11 \\ \hline
Quản lý phiên bản mã & Git (kèm Git LFS cho tệp lớn) & -- & 2.49.0 \\ \hline
\end{tabular}
\end{table}

\texttt{traci} và \texttt{sumolib} không cài từ PyPI mà nạp từ thư mục công cụ của bản cài SUMO, để khớp với tệp thực thi. \texttt{libsumo} và \texttt{sumo-data} cài không kèm phụ thuộc, vì phụ thuộc khai báo của chúng kéo theo \texttt{traci} và \texttt{sumolib} bản khác từ PyPI và che bản đi kèm SUMO. Một kiểm thử riêng xác nhận hai backend cho kết quả giống nhau (Mục~\ref{section:testing}).

LLM cục bộ là \texttt{qwen3.5:9b}: 9 tỷ tham số, lượng tử hóa Q4\_K\_M, tệp trọng số 5,6~GB, vừa với 8~GB bộ nhớ của GPU máy chính. Cấu hình vòng lặp: nhiệt độ 0,7; seed 42; 4 mẫu mỗi vòng; tối đa 3 lần thử mỗi mẫu; \texttt{max\_tokens} 7000; thinking tắt ở các chiến dịch ngày 08/11 (Mục~\ref{section:3.8}). Tùy chọn \texttt{--llm ollama} của vòng lặp chọn máy chủ Ollama mà không phải sửa tệp cấu hình LLM.

Trên máy chính, nếu \texttt{optuna} được nạp trước \texttt{libsumo} thì \texttt{libsumo} không nạp được: \texttt{optuna} nạp mô-đun \texttt{sqlite3} của Python, và tệp \texttt{sqlite3.dll} đi kèm Python xung đột với tệp cùng tên đi kèm \texttt{libsumo}. Nạp \texttt{libsumo} trước thì không lỗi, nên script Optuna nạp \texttt{libsumo} trước. Tệp nhật ký nghiên cứu của Optuna khóa bằng \texttt{JournalFileOpenLock}, vì cách khóa mặc định tạo liên kết tượng trưng (symlink), mà tài khoản Windows thường không có quyền tạo.

\subsection{Hạng mục đã hoàn thành}

Công việc chia thành 15 hạng mục (Bảng~\ref{table:4.4}).

\begin{table}[H]
\centering
\small
\caption{Các hạng mục và trạng thái.} \label{table:4.4}
\begin{tabular}{|p{0.8cm}|p{9.6cm}|p{2.6cm}|}
\hline
\textbf{\#} & \textbf{Sản phẩm} & \textbf{Trạng thái} \\ \hline
1 & Đọc cấu hình; cấu hình hiện trường đại diện có nguồn và mức bằng chứng cho từng giá trị & Hoàn thành \\ \hline
2 & Mạng mô phỏng sinh bằng script từ OSM và tệp vá; bảng cạnh; sơ đồ mạng & Hoàn thành \\ \hline
3 & vType B0, sinh nhu cầu, cấu hình SUMO và các đầu ra & Hoàn thành \\ \hline
4 & Trình chạy, manifest, chuyển đầu ra sang lược đồ quan sát & Hoàn thành \\ \hline
4b & Sửa bề rộng làn ra; tắt mặc định thiết bị đo chỉ số an toàn thay thế (SSM) của SUMO, gộp cảnh báo; chạy lại B0 trên dữ liệu đại diện & Hoàn thành \\ \hline
5 & Thí nghiệm khả thi về xe chờ giữa nút; quyết định dùng cơ chế sẵn có của SUMO & Hoàn thành \\ \hline
6 & Bảng giá trị mục tiêu và bộ sinh dữ liệu đại diện PROXY\_D1--PROXY\_D5 & Hoàn thành \\ \hline
7 & Miền tham số, lược đồ đặc tả, bộ kiểm tra hợp lệ, đặc tả B1 & Hoàn thành \\ \hline
8 & Bộ biên dịch đặc tả & Hoàn thành \\ \hline
9 & Bộ điều khiển TraCI và năm luật có điều kiện đầu tiên & Hoàn thành \\ \hline
10 & Chỉ số M1--M8, G0 và hàm khoảng cách & Hoàn thành \\ \hline
F2 & Mạng F2 (làn rẽ phải liên tục, \texttt{custom2}), thí nghiệm khả thi Phase 3A--3D, đoạn cọc, vỉa hè, luật \texttt{queue\_space\_seeking}, chỉ số M9--M10, đóng băng ngưỡng va chạm trên mạng F2 & Hoàn thành (đi ngược chiều hoãn); đoạn cọc đã bỏ ở mạng sau \\ \hline
KS & Cập nhật theo khảo sát hiện trường 28--30/10: đèn ba pha, hai nhánh cấm rẽ trái, hai làn vỉa hè, làn rẽ phải 30~m, kết nối từ lane 0 và 1, xét đèn theo liên kết của từng xe, prompt \texttt{v002}, luật \texttt{uturn\_bypass}, sinh lại dữ liệu đại diện & Hoàn thành; chiến dịch chuyển sang mạng 04/11 \\ \hline
QĐ & Quyết định 04/11 và 08/11: bỏ kết nối vi phạm từ lane 0, 1; luật vượt đỏ chờ nút thông; \texttt{uturn\_bypass} vào B1; lựa chọn cắt về; prompt \texttt{v003}--\texttt{v005}; M8 đọc làn ở đầu đoạn 30~m; P2 phải dùng đủ sáu luật & Hoàn thành \\ \hline
11 & Khung đánh giá chung và bộ đếm ngân sách & Hoàn thành \\ \hline
12 & Gọi LLM, prompt, phản hồi & Hoàn thành \\ \hline
13 & Vòng lặp LLM P1 (\texttt{params\_only}) và P2 (\texttt{full}) & Mã xong; P1, P2 với \texttt{qwen3.5:9b} đã chạy trên mạng 04/11 \\ \hline
14 & Phương án B2 và H bằng Optuna & Mã xong; B2 đã chạy trên bốn mạng, gồm mạng 04/11; H chưa chạy \\ \hline
15 & Đánh giá cuối và báo cáo & Dự kiến \\ \hline
\end{tabular}
\end{table}

\subsection{Kích thước mã nguồn}

Bảng~\ref{table:4.size} đếm tệp Python và số dòng (kể cả dòng trống và chú thích) của gói \texttt{thesis\_sumo} và bộ kiểm thử, tại ngày 08/11/2026.

\begin{table}[H]
\centering
\small
\caption{Kích thước mã nguồn theo gói.} \label{table:4.size}
\begin{tabular}{|l|r|r|}
\hline
\textbf{Gói hoặc mô-đun} & \textbf{Số tệp} & \textbf{Số dòng} \\ \hline
\texttt{runtime} & 12 & 1931 \\ \hline
\texttt{sim} & 4 & 1364 \\ \hline
\texttt{llm} & 4 & 1206 \\ \hline
\texttt{metrics} & 15 & 1158 \\ \hline
\texttt{net} & 5 & 1135 \\ \hline
\texttt{rules} & 4 & 968 \\ \hline
\texttt{proxy} & 2 & 910 \\ \hline
\texttt{baselines} & 2 & 373 \\ \hline
\texttt{demand} & 2 & 200 \\ \hline
Mô-đun cấp gốc (\texttt{config}, \texttt{evaluate}, \texttt{schema\_observed}, \texttt{\_\_init\_\_}) & 4 & 556 \\ \hline
\textbf{Tổng gói \texttt{thesis\_sumo}} & \textbf{54} & \textbf{9801} \\ \hline
Bộ kiểm thử (\texttt{tests/}) & 17 & 6948 \\ \hline
\end{tabular}
\end{table}

\section{Kiểm thử}
\label{section:testing}

Kiểm thử tự động dùng \texttt{pytest}. Phần lớn chạy trên dữ liệu nhỏ dựng tay hoặc trên mạng chữ thập tí hon, với giá trị kỳ vọng tính độc lập với mã được kiểm thử; một số chạy SUMO thật trong thời gian ngắn. Bảng~\ref{table:4.5} tóm tắt các bộ kiểm thử.

\begin{table}[H]
\centering
\small
\caption{Các bộ kiểm thử tự động.} \label{table:4.5}
\begin{tabular}{|p{3.6cm}|p{9.4cm}|}
\hline
\textbf{Bộ kiểm thử} & \textbf{Nội dung kiểm tra} \\ \hline
\texttt{test\_env} & \texttt{traci} và \texttt{sumolib} nạp từ bản cài SUMO và khớp phiên bản tệp thực thi \\ \hline
\texttt{test\_config} & Trường bắt buộc rỗng gây lỗi nêu tên trường; cấu hình đo thực không bao giờ tự lấy giá trị đại diện \\ \hline
\texttt{test\_net} & Dựng mạng bằng một lệnh; tệp vá, đèn, bảng cạnh; dựng lại cho kết quả giống từng byte; làn nội bộ có bề rộng bằng làn ra tương ứng \\ \hline
\texttt{test\_demand} & vType B0, tuyến ứng viên, sinh tuyến từ số đếm rẽ, kiểm tra dữ liệu số đếm \\ \hline
\texttt{test\_sumocfg} & Tùy chọn SUMO lấy từ cấu hình thí nghiệm, tệp đầu ra, chế độ SSM, gộp cảnh báo, chạy thật ngắn \\ \hline
\texttt{test\_runner} & Trình chạy, trường của manifest, chuyển sang lược đồ quan sát, tiêu chí đủ điều kiện \\ \hline
\texttt{test\_outputs\_\allowbreak synthetic} & Chuyển lược đồ trên đầu ra SUMO viết tay cho mạng chữ thập nhỏ, kết quả kỳ vọng tính tay \\ \hline
\texttt{test\_proxy} & Tính tất định của bộ sinh dữ liệu đại diện, lược đồ, nhất quán nội tại, độ khớp với giá trị mục tiêu \\ \hline
\texttt{test\_validator} & Miền tham số, tham số luật, mã lỗi và nội dung thông báo lỗi; lựa chọn cắt về; kiểm tra đủ luật của P2 \\ \hline
\texttt{test\_compiler} & Biên dịch B1 khớp từng byte với bản chụp; mọi thuộc tính nằm trong miền; SUMO nạp được với kiểm tra XSD; đặc tả B0 tương đương tệp vType viết tay; mã vType báo về khi lấy từ phân phối \\ \hline
\texttt{test\_runtime} & Mỗi luật một kịch bản trên mạng nhỏ; lệnh bị SUMO bỏ qua được ghi là không thực thi; khôi phục speed mode; tích hợp với trình chạy; giữ xe giữa nút không gây va chạm trong nút; cấm \texttt{ignoreIDs} trong mã \\ \hline
\texttt{test\_backends} & \texttt{libsumo} và \texttt{traci} cho quỹ đạo giống hệt nhau với cùng kịch bản và seed \\ \hline
\texttt{test\_metrics} & M1--M10, G0, hàm khoảng cách trên bảng dựng tay có đáp án; khoảng cách khi một phía rỗng không bao giờ nhỏ hơn khoảng cách khi sai \\ \hline
\texttt{test\_evaluate} & Ngân sách không vượt khi chạy lại; đặc tả hỏng không tốn lần chạy; mã trùng bị từ chối; chặn ngày đánh giá; quy tắc chọn đặc tả tốt nhất \\ \hline
\texttt{test\_llm} & Lược đồ cho LLM, khóa cấu trúc, phát lại nhật ký có kiểm định danh LLM, chặn rò rỉ dữ liệu ngày đánh giá \\ \hline
\texttt{test\_loop} & Chạy tiếp đúng sau khi bị ngắt (giữa vòng, ở đặc tả cuối còn đủ ngân sách) mà không trừ lại ngân sách; mẫu không hợp lệ không tốn lần chạy; khóa cấu trúc của P1; chạy thêm vòng khi còn ngân sách; số lần chạy đã dùng khớp bộ đếm \\ \hline
\texttt{test\_optuna} & Không gian tìm kiếm theo cấu trúc; mọi thử nghiệm ngẫu nhiên đều hợp lệ; chạy tiếp không trừ lại ngân sách \\ \hline
\end{tabular}
\end{table}

\textbf{Kết quả.} Bộ kiểm thử có 403 ca. Trên máy chính ngày 08/11/2026, 401 ca đạt và 2 ca thất bại dự kiến (xfail). Ngày 31/10 có 9 ca không đạt; nay chúng đều đạt. Sáu ca của \texttt{test\_optuna} đạt sau khi nhật ký Optuna đổi sang khóa \texttt{JournalFileOpenLock} (Mục~\ref{section:tools}), và kỳ vọng của \texttt{test\_junction\_info\_on\_proxy\_net} được cập nhật theo mạng 04/11, kèm phần kiểm tra lựa chọn cắt về. Cả hai xfail so sánh một tỷ lệ của ô tô trên năm ngày đại diện với giá trị mục tiêu, ngưỡng chấp nhận $\pm$15\%: (i) tỷ lệ xuất phát sớm của ô tô lệch $+49\%$; và (ii) tỷ lệ vượt đỏ của ô tô lệch $-21\%$.
Các độ lệch này đến từ số sự kiện nhỏ: cả năm ngày chỉ có khoảng 20 sự kiện ô tô vượt đỏ. Trên phiên bản kịch bản trước, 200 bộ sinh lại không cho thấy thiên lệch của tỷ lệ này. Bộ ngày không được sinh lại theo kết quả.

\textbf{Hai kiểm thử quan trọng nhất cho độ tin cậy.} Kiểm thử đối chiếu backend bảo đảm việc dùng \texttt{libsumo} để tăng tốc không đổi kết quả so với \texttt{traci}; ở cấp bộ điều khiển, hai backend cũng cho cùng nhật ký can thiệp. Kiểm thử giữ xe giữa nút chạy trên mạng chữ thập nhỏ có dòng cắt ngang và bật kiểm tra va chạm trong nút: giữ xe thuần bằng cách đặt tốc độ 0 không gây va chạm trong nút, còn một đối chứng âm (cho xe cắt ngang bỏ qua xe đang bị giữ) gây va chạm, nên kiểm thử phát hiện được nếu cách làm sai này quay lại. Script thí nghiệm vùng chờ giữa nút (Mục~\ref{section:box_capacity}) chưa có kiểm thử tự động.

\section{Thực nghiệm}
\label{section:experiments}

\subsection{Phạm vi và quy ước}
\label{subsection:net_versions}

Mọi thực nghiệm chạy trên mạng đại diện (cấu hình hiện trường lấy từ tài liệu và giả định). Kịch bản đổi nhiều lần theo khảo sát hiện trường, nên có sáu phiên bản mạng (Bảng~\ref{table:4.nets}). Số liệu giữa các mạng không so sánh trực tiếp được, vì đèn, nhu cầu và cấu hình đổi cùng lúc. Các chiến dịch từ ngày 08/11 dùng mạng 04/11.

\begin{table}[H]
\centering
\small\setlength{\tabcolsep}{3pt}
\caption{Các phiên bản mạng đại diện.} \label{table:4.nets}
\begin{tabular}{|p{2.0cm}|p{0.9cm}|p{2.3cm}|p{4.0cm}|p{2.2cm}|p{2.3cm}|}
\hline
\textbf{Mạng} & \textbf{Đèn} & \textbf{Cấm rẽ trái} & \textbf{Lane 0} & \textbf{Vỉa hè} & \textbf{Dùng ở} \\ \hline
Trước F2 & 2 pha & Lê Văn Lương tây & Chưa có làn rẽ phải liên tục & Không & Mục~\ref{section:box_capacity} đến~\ref{subsection:b1_check} \\ \hline
F2 (27/10) & 2 pha & Lê Văn Lương tây & Chỉ rẽ phải trên cả nhánh, cọc 30~m ngăn đổi làn & Nhánh tây & Mục~\ref{subsection:phase3} đến~\ref{subsection:interim} \\ \hline
28/10 đợt 1 & 3 pha & Hai nhánh Hoàng Đạo Thúy & Như F2 & Hai nhánh Lê Văn Lương & Mục~\ref{subsection:scenario_update} \\ \hline
28/10 đợt 2 & 3 pha & Hai nhánh Hoàng Đạo Thúy & Bỏ cọc, đi thẳng theo đèn & Hai nhánh Lê Văn Lương & Mục~\ref{subsection:scenario_update} \\ \hline
29/10 & 3 pha & Hai nhánh Hoàng Đạo Thúy & Làn rẽ phải chỉ ở 30~m cuối; thêm rẽ trái từ lane 0 và quay đầu từ lane 0, 1 cho xe vi phạm & Hai nhánh Lê Văn Lương & Mục~\ref{subsection:ref_n29} đến~\ref{subsection:uturn} \\ \hline
04/11 & 3 pha & Hai nhánh Hoàng Đạo Thúy & Làn rẽ phải chỉ ở 30~m cuối; rẽ trái và quay đầu chỉ từ lane 2 & Hai nhánh Lê Văn Lương & Mục~\ref{subsection:ref_n1104} đến~\ref{subsection:d1_backlog} \\ \hline
\end{tabular}
\end{table}

Với mỗi thực nghiệm, phần quan sát chỉ nêu số liệu, phần diễn giải chỉ dựa trên số liệu đó, và giả thuyết cần kiểm chứng nêu riêng.

Hai đại lượng dùng xuyên suốt. \emph{Va chạm/h} là số va chạm SUMO ghi nhận trong một lần chạy, quy về một giờ mô phỏng. \emph{Teleport kẹt} là số lần SUMO đưa một xe ra khỏi mạng vì xe kẹt quá thời gian cho phép (mặc định 300~s) hoặc kẹt khi nhường đường, không tính teleport do va chạm.

\subsection{Thí nghiệm khả thi về xe chờ giữa nút}
\label{section:box_capacity}

\textbf{Câu hỏi.} Với mô hình làn con SL2015, SUMO 1.27.0 có tự mô phỏng được xe máy rẽ trái trái phép dừng chờ giữa nút không, và chứa được bao nhiêu xe? Tiêu chí quyết định: nếu sức chứa tối đa của cơ chế sẵn có không nhỏ hơn trung bình mục tiêu của số xe chờ giữa nút (M7a, Bảng~\ref{table:3.3}), bằng 6 xe mỗi chu kỳ, thì dùng cơ chế sẵn có; nếu không, phải sửa mạng bằng một cạnh và tín hiệu riêng cho vùng chờ.

\textbf{Thiết lập.} Mạng đại diện, chu kỳ đèn 120~s (xanh, vàng, đỏ toàn phần 64/4/2~s cho Lê Văn Lương và 44/4/2~s cho Hoàng Đạo Thúy), vType B0. Nhu cầu tổng hợp riêng cho thí nghiệm, không phải dữ liệu đại diện. Mỗi nhánh Hoàng Đạo Thúy có 1200 xe máy/h và 250 ô tô/h đi thẳng. Dòng đi thẳng đối diện là dày (4000 xe máy/h và 700 ô tô/h, không còn khoảng trống trong pha xanh) hoặc vừa phải (1300 xe máy/h và 200 ô tô/h, gần lưu lượng đi thẳng của nhánh đối diện trong các ngày đại diện). Năm kịch bản khác nhau ở dòng xe vi phạm: (i) \texttt{sat}: 1500 xe/h, làm bão hòa vùng chờ; (ii) \texttt{target}: 225 xe/h, tức 7,5 xe vào mỗi chu kỳ theo trung bình mục tiêu; (iii) \texttt{sat\_tight}: 1500 xe/h, xe vi phạm có \texttt{minGap} 1,0~m và \texttt{minGapLat} 0,3~m; (iv) \texttt{sat\_mg18}: 1500 xe/h, \texttt{minGap} 1,8~m; và (v) \texttt{target\_modopp}: như \texttt{target}, với dòng đối diện vừa phải.
Mỗi kịch bản chạy 5 seed (1--5), mỗi lần 900~s (7 lần kết thúc pha xanh đối diện), tổng 25 lần chạy. Một bộ quan sát TraCI không can thiệp ghi vị trí, vị trí ngang và tốc độ của mọi xe trong vùng chờ, trạng thái đèn, teleport và va chạm ở mỗi bước 0,1~s. Xe dừng khi tốc độ dưới 5~km/h. \emph{Sức chứa} là số xe dừng lớn nhất trong vùng chờ khi dòng đối diện còn xanh. \emph{Thoát sớm} là xe rời vùng chờ qua khe trống trước khi pha xanh đối diện kết thúc, tính theo tỷ lệ trên số xe vào vùng chờ. Thí nghiệm chạy trước khi bật kiểm tra va chạm giữa các làn nội bộ cắt nhau.

\textbf{Quan sát (vùng chờ 3,2~m).} Bảng~\ref{table:4.box} cho kết quả theo kịch bản.

\begin{table}[H]
\centering
\small\setlength{\tabcolsep}{3pt}
\caption{Xe chờ giữa nút trên vùng chờ rộng 3,2~m (gộp chu kỳ 1--6 của 5 seed, 30 chu kỳ mỗi kịch bản). ``Chờ'': số xe đang chờ khi pha xanh đối diện kết thúc. ``Khoảng xả'': khoảng thời gian trung bình giữa hai xe liên tiếp rời vùng chờ.} \label{table:4.box}
\begin{tabular}{|p{2.3cm}|p{1.5cm}|p{1.3cm}|p{1.2cm}|p{0.9cm}|p{0.9cm}|p{2.4cm}|p{1.6cm}|}
\hline
\textbf{Kịch bản} & \textbf{Đối diện (xe/h)} & \textbf{minGap (m)} & \textbf{Chờ TB} & \textbf{p90} & \textbf{Max} & \textbf{Thoát sớm mỗi chu kỳ (tỷ lệ)} & \textbf{Khoảng xả (s)} \\ \hline
\texttt{sat} & 4700 & 2,5 & 8,00 & 8 & 8 & 0 (0\%) & 0,52 \\ \hline
\texttt{target} & 4700 & 2,5 & 6,43 & 8 & 8 & 0 (0\%) & 0,68 \\ \hline
\texttt{sat\_tight} & 4700 & 1,0 & 12,00 & 12 & 12 & 0 (0\%) & 0,51 \\ \hline
\texttt{target\_\allowbreak modopp} & 1500 & 2,5 & 2,60 & 6 & 7 & 4,17 (62\%) & 0,88 \\ \hline
\texttt{sat\_mg18} & 4700 & 1,8 & 10,00 & 10 & 10 & 0 (0\%) & 0,52 \\ \hline
\end{tabular}
\end{table}

Các điểm đáng chú ý: Trong mọi chu kỳ của mọi lần chạy, xe vi phạm dừng chờ trên làn nội bộ rẽ trái, xe đầu dừng cách điểm chờ nội bộ 0,10~m. Khi làn này đầy, các xe sau dừng tại vạch dừng. Với vType B0, sức chứa đo được là 8 xe ở cả 5 seed, xếp thành hai cột ngang ở vị trí ngang khoảng $-1{,}15$~m và từ 0 đến $+0{,}65$~m. Giảm \texttt{minGap} xuống 1,8~m và 1,0~m cho sức chứa 10 và 12 xe. Ở kịch bản \texttt{target}, vùng chờ đầy trong 43\% số chu kỳ. Xe chờ đầu tiên rời vùng chờ trung bình 4,7~s sau khi pha xanh đối diện kết thúc; vùng chờ trống sau 8,2~s với 8 xe và 10,4~s với 12 xe. Khi pha xanh Hoàng Đạo Thúy bắt đầu, còn 3,4 (\texttt{target}) đến 4,6 (\texttt{sat}) xe trong vùng chờ. Lượt qua vạch đầu tiên của làn 2 nhánh Hoàng Đạo Thúy phía bắc trễ 0,4--1,1~s với B0 và 1,6--2,7~s với 12 xe. Không có kẹt cứng và không có teleport không do va chạm trong cả 25 lần chạy; thời gian lưu lại vùng chờ lớn nhất là 70,4~s. Tổng cộng 334 va chạm trong 25 lần chạy; chỉ một va chạm có xe vi phạm tham gia, còn lại là giữa các xe đi thẳng cùng dòng.
Hình~\ref{fig:4.lateral} cho vị trí của các xe vi phạm đang dừng tại mỗi lần pha xanh đối diện kết thúc, gộp mọi seed của ba kịch bản bão hòa.

\begin{figure}[H]
\centering
\includegraphics[width=\textwidth]{Figure/box_lateral_32m.png}
\caption{Vị trí xe rẽ trái trái phép đang dừng trong vùng chờ rộng 3,2~m, mỗi lần pha xanh đối diện kết thúc, gộp 5 seed (kịch bản \texttt{sat}, \texttt{sat\_mg18}, \texttt{sat\_tight}). Trục hoành: vị trí đầu xe tính từ vạch dừng nhánh Lê Văn Lương phía tây dọc theo đường rẽ trái; trục tung: vị trí ngang. Hai đường ngang là mép làn nội bộ; đường đứt là điểm chờ nội bộ. Các xe xếp thành hai cột ngang.} \label{fig:4.lateral}
\end{figure}

\textbf{Diễn giải.} Cơ chế điểm chờ nội bộ sẵn có của SUMO tái hiện được cả chuỗi: vào nút, dừng chờ, chặn xe sau khi vùng chờ đầy, xả sau khi dòng đối diện dừng. Không cần tùy chọn hay sửa mạng. Sức chứa do hình học quyết định. Số hàng tính theo $\lfloor (18{,}79 + \text{minGap}) / (2{,}2 + \text{minGap}) \rfloor$, trong đó 18,79~m là vị trí đầu xe dẫn đầu; công thức dự đoán 4, 5 và 6 hàng tại \texttt{minGap} 2,5; 1,8 và 1,0~m, tức 8, 10 và 12 xe với hai cột, đúng như đo. Sức chứa 8 không nhỏ hơn trung bình mục tiêu 6, nên luận văn dùng cơ chế sẵn có.

Có hai điểm cần hiệu chỉnh. Sức chứa 8 nhỏ hơn phân vị 90 của mục tiêu (10 xe), nên phân phối số xe chờ bị cắt ở 8 với vType B0; đây là vấn đề tham số (\texttt{minGap} của nguyên mẫu \texttt{custom1}), không phải vấn đề mạng. Tỷ lệ thoát sớm phụ thuộc mạnh vào dòng đối diện: 0\% với dòng dày, 62\% với dòng gần mức ngày đại diện, trong khi mục tiêu là 20\% (giả định). Ở lưu lượng đại diện, tham số vì vậy chủ yếu phải giảm thoát sớm. Các va chạm của xe đi thẳng thuộc vấn đề bề rộng làn (Mục~\ref{subsection:exit_width}), không liên quan vùng chờ.

\textbf{Kiểm tra lại trên vùng chờ 3,5~m.} Sau khi bề rộng làn ra sửa thành 3,5~m (Mục~\ref{subsection:exit_width}), làn nội bộ rẽ trái cũng rộng 3,5~m. Kịch bản \texttt{target} được chạy lại với 5 seed, mỗi lần 900~s, vẫn chưa bật kiểm tra va chạm giữa các làn nội bộ cắt nhau (Bảng~\ref{table:4.box_r21}).

\begin{table}[H]
\centering
\small
\caption{Kiểm tra lại kịch bản \texttt{target} trên vùng chờ rộng 3,5~m (vType B0, 900~s mỗi seed).} \label{table:4.box_r21}
\begin{tabular}{|p{1.4cm}|p{3.4cm}|p{3.6cm}|p{2.3cm}|p{1.8cm}|}
\hline
\textbf{Seed} & \textbf{Số xe dừng lớn nhất trong vùng chờ} & \textbf{Số xe chờ TB khi hết xanh đối diện} & \textbf{Bề rộng dàn ngang (m)} & \textbf{Va chạm} \\ \hline
1 & 13 & 6,86 & 2,60 & 0 \\ \hline
2 & 8 & 6,00 & 2,60 & 0 \\ \hline
3 & 8 & 6,57 & 2,60 & 0 \\ \hline
4 & 7 & 4,86 & 2,50 & 1 \\ \hline
5 & 10 & 6,86 & 2,60 & 0 \\ \hline
\end{tabular}
\end{table}

Hình~\ref{fig:4.occupancy} cho diễn biến theo thời gian ở seed 1. Số xe dừng lớn nhất là 13, 8, 8, 7 và 10 theo seed 1--5, xếp thành ba cột ngang thay vì hai.

\begin{figure}[H]
\centering
\includegraphics[width=\textwidth]{Figure/box_occupancy_f2.png}
\caption{Số xe trong vùng chờ giữa nút rộng 3,5~m theo thời gian, seed 1. Trên: kịch bản \texttt{target} (dòng đối diện 4700 xe/h). Dưới: \texttt{target\_modopp} (dòng đối diện 1500 xe/h). Đường xanh: xe trên làn nội bộ; đường cam: xe đang dừng. Nền xám: pha xanh đi thẳng của nhánh Lê Văn Lương phía đông. Đường đứt: trung bình mục tiêu (6 xe); đường chấm: phân vị 90 (10 xe).} \label{fig:4.occupancy}
\end{figure} Số va chạm mỗi lần chạy là 0--1, so với 10--24 trên vùng chờ 3,2~m ở cùng kịch bản. M7a có trung bình 6,0, phân vị 90 bằng 10 và phân vị 95 bằng 11. Trên vùng chờ 3,5~m, sức chứa không còn là cận trên thấp hơn mục tiêu: giá trị 13 ở seed 1 đã vượt phân vị 95. Quyết định dùng cơ chế sẵn có giữ nguyên. Kết luận về sức chứa và tỷ lệ thoát sớm ở trên đo trên vùng chờ 3,2~m; riêng tỷ lệ thoát sớm chưa đo lại trên vùng chờ 3,5~m với nhiều seed.

\textbf{Giả thuyết cần kiểm chứng.} Sức chứa phụ thuộc trực tiếp vào hình học vùng chờ (chiều dài làn nội bộ, vị trí điểm chờ, bề rộng làn), hiện lấy từ tiêu chuẩn thiết kế chứ không đo. Khi có hình học đo thực, thí nghiệm cần chạy lại; nếu sức chứa khi đó nhỏ hơn mục tiêu, phương án sửa mạng cần xét lại.

\subsection{Ảnh hưởng của bề rộng làn ra đến va chạm}
\label{subsection:exit_width}

\textbf{Thiết lập.} Trong mạng đầu tiên, cấu hình đại diện chỉ đặt bề rộng 3,5~m cho làn vào; làn ra giữ mặc định 3,2~m của \texttt{netconvert}, không có nguồn. B0 chạy trên dữ liệu SYNTHETIC (chỉ để thử đường ống) và trên PROXY\_D1, mỗi trường hợp 3 seed, chưa bật kiểm tra va chạm giữa các làn nội bộ cắt nhau. Lần chạy PROXY\_D1 dùng trực tiếp tệp thực thi SUMO, không ghi quỹ đạo, SSM và detector, vì trình chạy khi đó dừng giữa chừng do SSM ghi quá nhiều cảnh báo.

\textbf{Quan sát.} Va chạm/h theo seed 1, 2, 3 là 5, 16, 10 trên SYNTHETIC và 62, 52, 46 trên PROXY\_D1. Trong 31 va chạm của các lần chạy SYNTHETIC, 20 xảy ra khi xe vào nút, 8 ở cuối làn gần vạch dừng, 3 là va chạm làn con khi chuyển làn cách vạch dừng 16--18~m. Làn ra rộng 3,5~m giảm tổng số va chạm từ 31 xuống 2 trên SYNTHETIC và từ 160 xuống 10 trên PROXY\_D1. Trên PROXY\_D1, số teleport kẹt khi đó là 18, 32 và 10 theo seed. Thử từng yếu tố khác: độ phân giải làn con 0,2~m cho 40 va chạm, 0,8~m cho 12 va chạm; bật kiểm tra va chạm trong nút làm số va chạm tăng từ 31 lên 86.

\textbf{Diễn giải.} \texttt{netconvert} gán cho làn nội bộ bề rộng của làn ra. Xe máy trong mô hình làn con đi sát vạch làn được (vị trí ngang tới $\pm$1,5~m). Khi xe đi từ làn vào 3,5~m sang làn nội bộ 3,2~m mà tâm các làn vẫn cách nhau 3,5~m, các xe máy đi song song chồng lấn nhau. Đây là hệ quả của một giá trị mặc định do chính công cụ để lại, không phải hành vi giao thông. Hình~\ref{fig:4.gui} là một trường hợp như vậy.

\begin{figure}[H]
\centering
\includegraphics[width=0.75\textwidth]{Figure/gui_collision_32m.png}
\caption{Ảnh \texttt{sumo-gui} trên mạng ban đầu (làn ra 3,2~m), seed 2, $t = 2647$~s: hai xe máy đi song song ở làn 1 và làn 2 sắp vào nút. 4,7~s sau, tại $t = 2651{,}7$~s, SUMO ghi một va chạm trong nút giữa hai xe này (khoảng hở ngang $-0{,}26$~m).} \label{fig:4.gui}
\end{figure}

Cách sửa: khai báo danh sách làn ra trong cấu hình hiện trường với bề rộng bằng làn vào (3,5~m, theo tiêu chuẩn thiết kế đường đô thị~\cite{tcvn13592}), kèm một kiểm thử khẳng định mọi làn nội bộ có bề rộng bằng làn ra tương ứng. Sau khi sửa, cả 56 làn của mạng rộng 3,5~m. Lần chạy kiểm tra 600~s trên PROXY\_D1, seed 1, tắt SSM, cho 0 va chạm và 0 teleport trên mạng mới, so với 8 va chạm và 8 teleport (đều do va chạm) trên mạng cũ.

\textbf{Giả thuyết cần kiểm chứng.} Số va chạm của mô hình làn con rất nhạy với bề rộng làn nội bộ. Nếu mạng đo thực có làn vào và làn ra khác nhau, va chạm kiểu này có thể xuất hiện lại; khi đó phải báo cáo như một giới hạn của mô hình làn con, không coi là hành vi thực.

\subsection{Cấu hình cơ sở B0 trên dữ liệu đại diện}
\label{subsection:b0_proxy}

\textbf{Thiết lập.} Sau khi sửa bề rộng làn ra, B0 chạy trên ba ngày hiệu chỉnh PROXY\_D1--PROXY\_D3, mỗi ngày 3 seed (1, 2, 3), mỗi lần 1800~s (thay cho 3600~s, vì một lần PROXY\_D1 3600~s mất hơn 20 phút), tắt SSM. Các lần chạy này diễn ra \emph{trước} khi luận văn quyết định bật kiểm tra va chạm trong nút cho mọi lần chạy, nên số va chạm không gồm va chạm giữa các làn nội bộ cắt nhau.

Bảng~\ref{table:4.b0} dùng bốn đại lượng: teleport kẹt (định nghĩa ở Mục~\ref{section:experiments}); số xe chưa chèn được vào mạng khi kết thúc; cờ \texttt{valid\_strict} (không có teleport và va chạm nào); tỷ lệ ô có GEH $<5$ của cổng G0 (Bảng~\ref{table:3.3}) tính trên 62 ô theo 15 phút $\times$ nhánh $\times$ hướng $\times$ loại xe, và tỷ lệ tương ứng khi gộp loại xe (ô theo nhánh $\times$ hướng).

\begin{table}[H]
\centering
\small\setlength{\tabcolsep}{3pt}
\caption{B0 trên PROXY\_D1--PROXY\_D3 (1800~s, chưa bật kiểm tra va chạm trong nút, tắt SSM).} \label{table:4.b0}
\begin{tabular}{|p{1.3cm}|p{0.9cm}|p{1.5cm}|p{1.6cm}|p{1.8cm}|p{1.8cm}|p{1.6cm}|p{1.7cm}|}
\hline
\textbf{Ngày} & \textbf{Seed} & \textbf{Va chạm/h} & \textbf{Teleport kẹt} & \textbf{Xe chưa chèn cuối} & \textbf{valid\_\allowbreak strict} & \textbf{GEH$<$5 (ô)} & \textbf{GEH$<$5 (nhánh $\times$ hướng)} \\ \hline
D1 & 1 & 2 & 14 & 161 & Không & 0,919 & 0,844 \\ \hline
D1 & 2 & 0 & 22 & 0 & Không & 0,952 & 0,875 \\ \hline
D1 & 3 & 2 & 10 & 0 & Không & 1,000 & 0,969 \\ \hline
D2 & 1 & 2 & 22 & 0 & Không & 1,000 & 1,000 \\ \hline
D2 & 2 & 6 & 23 & 0 & Không & 1,000 & 0,969 \\ \hline
D2 & 3 & 0 & 2 & 0 & Không & 1,000 & 1,000 \\ \hline
D3 & 1 & 0 & 0 & 0 & Có & 1,000 & 1,000 \\ \hline
D3 & 2 & 2 & 11 & 0 & Không & 0,984 & 0,906 \\ \hline
D3 & 3 & 0 & 0 & 0 & Có & 1,000 & 1,000 \\ \hline
\end{tabular}
\end{table}

\textbf{Quan sát.} Va chạm/h của 9 lần chạy là 2, 0, 2, 2, 6, 0, 0, 2, 0 (trung bình 1,56); với thời lượng 1800~s, mỗi va chạm tương ứng 2 va chạm/h. So với 46--62 va chạm/h trên mạng làn ra 3,2~m, mức này giảm mạnh. 7/9 lần chạy có teleport kẹt; chỉ hai lần (PROXY\_D3 seed 1 và 3) đạt \texttt{valid\_strict}; lần PROXY\_D1 seed 1 kết thúc với 161 xe chưa chèn được. Mọi teleport kẹt đều là xe máy, xảy ra thành cụm trong cùng giây; lần đầu hầu như luôn ở vạch dừng làn 2 (làn chung cho đi thẳng, rẽ trái và rẽ phải) của nhánh Lê Văn Lương phía tây, rồi lan sang các liên kết làn 2 trong nút. PROXY\_D1 có nhu cầu cao hơn PROXY\_D2 và PROXY\_D3 khoảng 20\% (4556, 3765 và 3857 xe trong 1800~s). GEH $<5$ theo loại xe nằm trong 0,919--1,000, đều vượt ngưỡng 0,85. Ở mức nhánh $\times$ hướng, tỷ lệ là 0,844--1,000; giá trị thấp nhất thuộc lần chạy còn xe chưa chèn.

\textbf{Thí nghiệm một yếu tố.} Trên PROXY\_D1 với 3 seed, bỏ 217 xe rẽ trái trái phép (2,4\% số xe) làm tổng số teleport còn 0, 0 và 3. Bỏ mọi xe rẽ trái cho 0, 3 và 2 teleport, không còn teleport kẹt. Giảm \texttt{minGap} xuống 1,5~m làm kết quả xấu hơn (68 và 103 teleport ở seed 1 và 2).

\textbf{Diễn giải.} Nguyên nhân chính gây kẹt của B0 trên dữ liệu đại diện là xe máy rẽ trái trái phép chờ giữa nút với vType mặc định, làm tắc làn 2 của nhánh có xe vi phạm. Đây là một phát hiện về hành vi của cấu hình mặc định, không phải lỗi mạng, nên mạng không được vá thêm; giảm kẹt này là mục tiêu mà các bộ luật hiệu chỉnh phải đạt. GEH cao chỉ cho thấy nút phục vụ được nhu cầu đại diện: số ``quan sát'' đại diện và tệp tuyến cùng sinh từ một bộ sinh, nên G0 ở đây không chứng minh độ thực tế.

\textbf{Giới hạn về điều kiện chạy.} Từ 9 tỷ lệ va chạm trên, ngưỡng va chạm từng được đặt bằng $\max(2, \lceil p_{90} \rceil) = 3$ va chạm/h (với $p_{90} = 2{,}8$). Ngưỡng này \emph{đã bị hủy} khi luận văn quyết định bật kiểm tra va chạm trong nút cho mọi lần chạy: khi tắt, va chạm giữa xe chờ giữa nút và dòng cắt ngang không được đếm. Hai ngưỡng mới (va chạm trên làn và trong nút) được tính lại từ B0 với kiểm tra bật trên mạng F2 (Mục~\ref{subsection:r28_f2}). Vì vậy số va chạm ở Bảng~\ref{table:4.b0} không dùng để so với các cấu hình chạy sau quyết định này. Kết quả mục này chỉ phản ánh hành vi mô hình trên dữ liệu đại diện và chỉ phủ nửa giờ đầu của khung cao điểm.

\subsection{Kiểm tra chức năng bộ luật thủ công B1}
\label{subsection:b1_check}

\textbf{Thiết lập.} Đặc tả B1 (tham số vType chuyển từ Saval và cộng sự~\cite{saval2023mixed}; luật vượt đèn đỏ theo đoàn với $p_0 = p_{\max} = 0{,}5$, $k = 0$; luật xuất phát sớm với $p = 0$) chạy một lần 600~s trên PROXY\_D1, seed 1, bật kiểm tra va chạm trong nút, trên máy đang chạy song song các tác vụ khác. Lần chạy này kiểm tra bộ điều khiển hoạt động đúng trong trình chạy; nó không phải một đánh giá, và chưa có lần B0 cùng điều kiện để so sánh.

\textbf{Quan sát.} Chỉ dùng vType B1, và dùng thêm bộ điều khiển với các luật không bao giờ kích hoạt: hai lần chạy cho cùng kết quả (0 va chạm trên làn, 7 va chạm trong nút, 11 teleport, trong đó 4 không do va chạm). Quỹ đạo giống hệt nhau: 628 bản ghi chuyến đi, không có khác biệt. Với đầy đủ B1: 0 va chạm trên làn, 13 va chạm trong nút (1 liên quan xe đang chịu can thiệp), 13 teleport, không teleport nào không do va chạm. Luật vượt đỏ theo đoàn đưa ra 80 quyết định; 53 được SUMO thực thi, 27 không thực thi vì pha đỏ đã hết trước khi xe qua vạch. Luật xuất phát sớm với $p = 0$ không tạo can thiệp nào qua 69 lần rút ngẫu nhiên. Thời gian xử lý của bộ điều khiển là 3--7~s cho mỗi 600~s mô phỏng.

\textbf{Diễn giải.} Bộ điều khiển không đổi quỹ đạo khi không có luật nào kích hoạt, ghi đúng các lệnh bị SUMO bỏ qua, và có chi phí tính toán nhỏ. Số va chạm và teleport đến từ một seed, 600~s, nên không dùng để kết luận về chất lượng của B1.

\subsection{Thí nghiệm khả thi Phase 3A--3D trên mạng có làn rẽ phải liên tục}
\label{subsection:phase3}

\textbf{Bối cảnh.} Khi xác nhận mọi nhánh có làn rẽ phải liên tục, mạng F2 coi lane 0 của mỗi nhánh là làn chỉ để rẽ phải, có cọc ngăn đổi làn (không thêm làn mới), nên mỗi nhánh còn hai làn đi thẳng có đèn. Giả định về cọc và về cấm đi thẳng về sau được thay (cuối mục này). Trước khi cài hành vi 6--8 và xác nhận lại hành vi 10 trên mạng này, luận văn chạy bốn thí nghiệm khả thi. Ba thí nghiệm 3A, 3B, 3C chạy trên kịch bản đồ chơi riêng, nên các con số trong đó là thiết lập thí nghiệm, không phải số đo. Thí nghiệm 3D chạy trên mạng đại diện F2. Mọi lần chạy bật kiểm tra va chạm trong nút.

\textbf{Quan sát.} Bảng~\ref{table:4.phase3} tóm tắt bốn thí nghiệm.

\begin{table}[H]
\centering
\small
\caption{Kết quả bốn thí nghiệm khả thi Phase 3A--3D.} \label{table:4.phase3}
\begin{tabular}{|p{2.4cm}|p{11.6cm}|}
\hline
\textbf{Thí nghiệm} & \textbf{Kết quả} \\ \hline
3A: chặn làn rẽ phải (hành vi 6) & SL2015 không tự tạo hành vi này: trong 12 lần chạy với hàng chờ tới 232~m, không xe vi phạm nào vào lane 0, vì với tuyến đi thẳng lane 0 là ngõ cụt. Hành vi chỉ xuất hiện khi có điều kiện kích hoạt qua TraCI theo độ dài hàng chờ, kết hợp tách đoạn có cọc thành cạnh riêng cấm chuyển làn (giá trị \texttt{ignoring} của thuộc tính cấm chuyển làn bị \texttt{netconvert} 1.27 bỏ qua không báo, nên dùng \texttt{emergency}). Khi đó xe vi phạm đi trên lane 0 và dừng trước đầu đoạn cọc để chờ nhập lại. Xe rẽ phải trễ 4,2~s với xác suất kích hoạt 0,3 và 8--27~s với xác suất 1, so với 3,4--3,6~s khi không có xe chặn. Không có va chạm trong nút trong 90 lần chạy; va chạm trên làn tăng theo xác suất kích hoạt (0--4 mỗi lần chạy); xe tuân thủ không bị kẹt \\ \hline
3B: đi trên vỉa hè (hành vi 7) & SL2015 không tự dùng làn vỉa hè (nhiều nhất 1 trên 134 xe). Trong các cấu trúc mạng đã thử, chỉ cấu trúc ``cạnh vỉa hè riêng chạy tới vạch dừng, có liên kết đèn riêng và được ưu tiên'' kèm điều kiện kích hoạt là chạy được: không có teleport kẹt, ít nhất 97\% xe đi hết vỉa hè; chi phí cho xe tuân thủ là thời gian đi thẳng tăng từ 42~s lên 59--84~s. Cấu trúc nhập lại kiểu phải nhường đường gây kẹt cứng; thêm làn mép trên cùng cạnh làm dịch mọi chỉ số làn nên không dùng. SUMO không cung cấp chỉ số an toàn cho xe trên vỉa hè \\ \hline
3C: đi ngược chiều (hành vi 8) & Cơ chế vượt bằng làn ngược chiều sẵn có không dùng được: trên đường nhiều làn chỉ xe ưu tiên được dùng; tham số \texttt{lcOpposite} theo tài liệu chỉ áp dụng cho mô hình LC2013, không chạy cùng mô hình làn con; với SL2015 và làn con, xe có phản ứng nhưng đứng lại trên làn ngược chiều (trung vị 54~s). Một làn ảo ngược chiều dành cho \texttt{custom1} chạy được (71/71 xe dùng, 0 va chạm) nhưng SUMO không tính xung đột giữa làn ảo và dòng ngược chiều thật, nên năng lực bị thổi phồng và không có chỉ số an toàn. Lệnh TraCI \texttt{changeSublane} đưa được xe sang làn ngược chiều thật (có tính xung đột) nhưng không kiểm soát được điểm quay về \\ \hline
3D: chờ giữa nút (hành vi 10) & Trên mạng F2, kịch bản \texttt{target}, 5 seed: số xe dừng lớn nhất là 13, 8, 8, 7 và 10, giống kết quả trên vùng chờ 3,5~m (Bảng~\ref{table:4.box_r21}). Không có va chạm giữa xe rẽ trái trái phép với xe khác và không có kẹt cứng \\ \hline
\end{tabular}
\end{table}

\textbf{Quyết định.} Hành vi 6 và 7 được cài vào mạng F2 với giá trị giả định lấy từ thiết lập của 3A và 3B: đoạn cọc 30~m; vỉa hè ở nhánh Lê Văn Lương phía tây, rộng 3,0~m, tốc độ 5~m/s, bắt đầu 170~m trước vạch dừng; kèm lựa chọn \texttt{right\_lane} và \texttt{sidewalk} của luật \texttt{queue\_space\_seeking}. Trong tệp cấu hình hiện trường đo thực, các trường này để trống. Hành vi 8 được \emph{hoãn}: không cài vào mô hình, chỉ số tương ứng (M11) và lựa chọn \texttt{opposite} của luật không được hiện thực, và hành vi này nằm trong phần hạn chế. Hành vi 10 tiếp tục dùng cơ chế điểm chờ nội bộ sẵn có.

\textbf{Kiểm tra chức năng sau khi cài.} Trên PROXY\_D1, 3 seed, 1800~s, so sánh B1 với mọi xác suất bằng 0 và một biến thể có xác suất chọn làn rẽ phải 0,3, chọn vỉa hè 0,3, $q = 30$~m. Lựa chọn \texttt{right\_lane} được SUMO thực thi ở 93/100, 132/143 và 104/109 lần quyết định. Lựa chọn \texttt{sidewalk} chỉ thực thi ở 8/17, 16/33 và 3/25 lần, chủ yếu vì xe ở làn 1--2 không kịp sang lane 0 trong khoảng 27~m trước điểm bắt đầu vỉa hè. Va chạm trong nút không có xu hướng theo luật; số va chạm liên quan can thiệp là 3--11, so với 1--5 của B1.

\textbf{Giả thuyết cần kiểm chứng.} Mức trễ của xe rẽ phải và tỷ lệ dùng vỉa hè phụ thuộc trực tiếp vào chiều dài đoạn làn rẽ phải, vị trí điểm lên vỉa hè và điểm nhập lại, đều chưa đo. Cần đo các đại lượng này tại nút trước khi diễn giải M9, M10.

\textbf{Phần không còn áp dụng trên mạng 29/10.} Khảo sát hiện trường sau các thí nghiệm này cho thấy không có cọc, lane 0 cho đi thẳng khi đèn xanh, và làn rẽ phải chỉ dài 30~m. Kết luận của 3A rằng SUMO không tự tạo hành vi 6 vì vậy chỉ đúng khi lane 0 là ngõ cụt với tuyến đi thẳng; trên mạng 29/10 hành vi này xuất hiện ở cả B0 (Mục~\ref{subsection:ref_n29}), và lựa chọn \texttt{right\_lane} không còn tồn tại. Các con số về độ trễ, về \texttt{ignoring} và \texttt{emergency} của 3A là kết quả của cấu trúc có cọc. Kết quả 3B và 3C giữ nguyên. Kết quả 3D đo trên vùng chờ của nhánh Lê Văn Lương phía tây; sức chứa hai vùng chờ Hoàng Đạo Thúy chưa đo.

\subsection{Ngưỡng va chạm đóng băng trên mạng F2}
\label{subsection:r28_f2}

\textbf{Thiết lập.} Trước vòng LLM đầu tiên, B0 chạy trên mạng F2 với PROXY\_D1--PROXY\_D3, mỗi ngày 3 seed, 1800~s, bật kiểm tra va chạm trong nút, tắt SSM. Hai ngưỡng được tính bằng $\max(2, \lceil p_{90} \rceil)$ của 9 tỷ lệ va chạm mỗi giờ, đóng băng trong cấu hình thí nghiệm và không đổi sau khi thấy kết quả các vòng tối ưu (Bảng~\ref{table:4.r28}).

\begin{table}[H]
\centering
\small
\caption{B0 trên mạng F2, PROXY\_D1--PROXY\_D3 (1800~s mỗi lần, bật kiểm tra va chạm trong nút); số va chạm là số đếm trong một lần chạy.} \label{table:4.r28}
\begin{tabular}{|p{1.6cm}|p{1.8cm}|p{2.0cm}|p{3.4cm}|p{2.2cm}|}
\hline
\textbf{Ngày / seed} & \textbf{Va chạm làn} & \textbf{Va chạm trong nút} & \textbf{Teleport (tổng / không do va chạm)} & \textbf{Xe chưa chèn cuối} \\ \hline
D1 / 1 & 2 & 41 & 59 / 16 & 339 \\ \hline
D1 / 2 & 0 & 53 & 79 / 26 & 101 \\ \hline
D1 / 3 & 2 & 54 & 56 / 0 & 0 \\ \hline
D2 / 1 & 0 & 47 & 47 / 0 & 0 \\ \hline
D2 / 2 & 0 & 69 & 69 / 0 & 0 \\ \hline
D2 / 3 & 1 & 63 & 91 / 27 & 51 \\ \hline
D3 / 1 & 1 & 78 & 79 / 0 & 0 \\ \hline
D3 / 2 & 0 & 63 & 78 / 15 & 0 \\ \hline
D3 / 3 & 0 & 53 & 60 / 7 & 0 \\ \hline
\end{tabular}
\end{table}

\textbf{Quan sát.} Quy về mỗi giờ, va chạm trên làn là 4, 0, 4, 0, 0, 2, 2, 0, 0 (trung bình 1,333; $p_{90} = 4{,}0$), còn va chạm trong nút nằm trong 82--156 ($p_{90} = 141{,}6$). Các ngưỡng đóng băng là \textbf{4 va chạm trên làn/h} và \textbf{142 va chạm trong nút/h}; cổng tạm thời ``B0 không quá 5 va chạm trên làn/h'' đạt. Không lần chạy B0 nào đủ điều kiện: các lần bị loại vì teleport không do va chạm, xe chưa chèn được khi kết thúc, hoặc va chạm trong nút.

\textbf{Diễn giải.} Với hai làn đi thẳng mỗi nhánh, mạng đại diện quá tải ở mức nhu cầu đại diện 8180 xe/h. Trong kiểm tra chức năng nhóm F2, số xe chưa chèn được cuối lần chạy là 300--1200 ở cả B0 và B1. Luận văn giữ nguyên nhu cầu và yêu cầu không còn xe chưa chèn khi kết thúc, không sửa mạng để giảm tải. Khi không có đặc tả nào đủ điều kiện, đặc tả tốt nhất được chọn theo số teleport kẹt mỗi giờ rồi theo $S$ (Mục~\ref{section:3.8}). Đây là giới hạn của mô hình nền và phải nêu cùng mọi kết quả hiệu chỉnh. Ngưỡng 142 va chạm trong nút/h gần như không loại lần chạy nào; nó chủ yếu ghi nhận rằng va chạm trong nút của mô hình làn con rất phổ biến ngay ở cấu hình mặc định.

\subsection{Kết quả tạm thời của các chiến dịch hiệu chỉnh trên mạng F2}
\label{subsection:interim}

\textbf{Phạm vi và điều kiện.} Mục này báo cáo các chiến dịch trên mạng F2 tại 27/10/2026. Kết quả trên các mạng sau nằm ở Mục~\ref{subsection:scenario_update} và~\ref{subsection:ref_n29}.
B2 (Optuna trên cấu trúc B1) đã dùng hết 36/36 lần chạy, tức 12 thử nghiệm. Để có mốc, B0 và B1 được chấm theo cùng cách trong chiến dịch tham chiếu REF. P1 và P2 với LLM cục bộ \texttt{qwen3.5:9b} chưa có kết quả chính thức, nên chưa có mặt ở đây (Mục~\ref{subsection:planned_experiments}). Mỗi đặc tả được đánh giá bằng 3 lần chạy 1800~s (PROXY\_D1, D2, D3, seed 1), trên dữ liệu \texttt{LITERATURE\_PROXY}.
Các tỷ lệ mỗi giờ là trung bình của 3 lần chạy; vì mỗi lần dài 0,5~h, một sự kiện tương ứng 2 sự kiện/h. Bảng~\ref{table:4.interim_summary} tóm tắt theo chiến dịch; Bảng~\ref{table:4.interim_specs} ở Mục~\ref{section:appendix_specs} liệt kê từng đặc tả.

\begin{table}[H]
\centering
\small\setlength{\tabcolsep}{3pt}
\caption{Tóm tắt theo chiến dịch (PROXY\_D1--D3, seed 1, 1800~s). Mọi đặc tả đều bị loại bởi tiêu chí đủ điều kiện. ``Teleport kẹt/h'' gồm teleport do kẹt và do kẹt khi nhường đường. ``Tốt nhất'' theo quy tắc dự phòng: ít teleport kẹt/h nhất, rồi $S$ nhỏ nhất.} \label{table:4.interim_summary}
\begin{tabular}{|p{1.5cm}|p{1.5cm}|p{2.6cm}|p{1.5cm}|p{1.7cm}|p{1.5cm}|p{2.6cm}|}
\hline
\textbf{Chiến dịch} & \textbf{Đặc tả (lần chạy)} & \textbf{$S$ nhỏ nhất / trung vị / lớn nhất} & \textbf{Va chạm làn/h} & \textbf{Va chạm tổng/h} & \textbf{Teleport kẹt/h} & \textbf{Tốt nhất (quy tắc dự phòng)} \\ \hline
B2 & 12 (36) & 1,182 / 1,329 / 1,523 & 2,0--46,0 & 90,7--168,0 & 12,0--76,7 & \texttt{B2\_\allowbreak t010}: 12,0/h; $S$ 1,374 \\ \hline
B0 (REF) & 1 (3) & 1,181 & 2,0 & 112,7 & 10,7 & -- \\ \hline
B1 (REF) & 1 (3) & 1,216 & 0,0 & 107,3 & 45,3 & -- \\ \hline
\end{tabular}
\end{table}

\textbf{Quan sát.} Không đặc tả nào của B2 đủ điều kiện. Mọi đặc tả có teleport kẹt (ít nhất 12,0/h), trong khi tiêu chí yêu cầu 0; đặc tả tốt nhất vì vậy chọn bằng quy tắc dự phòng. B0 và B1 chấm theo cùng cách cũng bị loại. Đặc tả tốt nhất theo quy tắc dự phòng là \texttt{B2\_t010} (12,0 teleport kẹt/h), dù nó có va chạm trên làn cao nhất chiến dịch (46,0/h) và $S$ (1,374) cao hơn đặc tả có $S$ nhỏ nhất (\texttt{B2\_t007}, 1,182); quy tắc dự phòng không áp mức trần va chạm. $S$ tốt nhất của B2 giảm từ 1,302 ở thử nghiệm đầu xuống 1,182 ở thử nghiệm thứ 8 và không giảm thêm trong 4 thử nghiệm cuối. So với B0 ($S = 1{,}181$; 10,7 teleport kẹt/h): không đặc tả B2 nào có $S$ nhỏ hơn, cũng không đặc tả nào có ít teleport kẹt hơn. So với B1, cấu trúc xuất phát của B2 ($S = 1{,}216$; 0,0 va chạm trên làn/h; 45,3 teleport kẹt/h): 1/12 đặc tả có $S$ nhỏ hơn, 8/12 đặc tả có ít teleport kẹt hơn.

\textbf{Diễn giải có cơ sở.} Trên mạng đại diện quá tải, quy tắc chọn dự phòng ưu tiên giảm ùn tắc hơn độ tương đồng, nên ``đặc tả tốt nhất'' có thể có $S$ cao hơn hoặc va chạm nhiều hơn các đặc tả khác cùng chiến dịch, như \texttt{B2\_t010}. Trong ngân sách 36 lần chạy, Optuna chưa tìm được bộ tham số nào trên cấu trúc B1 tốt hơn cấu hình mặc định B0 theo $S$ hay theo số teleport kẹt.

\textbf{Giới hạn của kết quả.} Bốn điều kiện hạn chế cách đọc các số trên: (i) mỗi ngày chỉ một seed, nên chưa tách được biến động do seed; (ii) mọi đặc tả bị loại bởi tiêu chí đủ điều kiện, và đặc tả tốt nhất chọn theo quy tắc dự phòng; (iii) mô hình nền quá tải ở nhu cầu đại diện; và (iv) dữ liệu mục tiêu là dữ liệu đại diện, và các khóa có mục tiêu giả định chiếm hơn một nửa giá trị $S$ ở mọi đặc tả.
Đánh giá E1, E2 trên ngày đánh giá với nhiều seed và kiểm định thống kê (Mục~\ref{subsection:planned_experiments}) là bước cần có trước khi kết luận so sánh.

\subsection{Cập nhật kịch bản và kết quả trên hai mạng ngày 28/10}
\label{subsection:scenario_update}

\textbf{Bối cảnh.} Khảo sát hiện trường ngày 28/10/2026 cho thấy bốn điểm khác mạng F2: (i) đèn có ba pha, rẽ trái từ Lê Văn Lương có mũi tên riêng; (ii) chỉ hai nhánh Hoàng Đạo Thúy cấm rẽ trái; (iii) xe máy đi vỉa hè ở cả hai chiều Lê Văn Lương; và (iv) không có cọc, và đi thẳng từ lane 0 khi đèn xanh là đúng luật.
Ba điểm đầu vào mạng 28/10 đợt 1, điểm cuối vào đợt 2. Mạng, nhu cầu, năm ngày đại diện và prompt (\texttt{v002}) được dựng lại. Bộ điều khiển và bộ chuyển đầu ra đổi sang xét đèn theo liên kết của từng xe, vì cách cũ tính xe rẽ trái Lê Văn Lương trong pha mũi tên là vượt đỏ.

\textbf{Hai lỗi mạng phát hiện khi chạy thử.} Lỗi thứ nhất: xe chờ giữa nút ở Hoàng Đạo Thúy bị xe đi thẳng đối diện đâm, vì \texttt{netconvert} đặt điểm chờ với giả định xe rộng 1,8~m, trong khi xe máy chờ sát mép làn nội bộ lấn sang dòng đi thẳng đang xanh cùng pha. Sửa bằng lượt \texttt{netconvert} cuối (Bảng~\ref{table:3.0}); trên PROXY\_D1, seed 1, 2, 3, số va chạm của B0 từ 32, 22, 48 còn 11, 29, 17. Tác dụng phụ: rẽ trái nhánh Lê Văn Lương phía đông nay nhường rẽ trái nhánh phía tây, và độ trễ tương ứng chưa đo. Lỗi thứ hai: kẹt cứng ở cửa ra vỉa hè, khi dòng rẽ phải từ vỉa hè và dòng đi thẳng Hoàng Đạo Thúy cắt ngang nhường nhau cả hai chiều. Sửa bằng một \texttt{prohibition} chỉ giữ chiều xe vỉa hè nhường.

\textbf{Thiết lập.} B0 và B1 chấm cùng cách với Mục~\ref{subsection:interim} (PROXY\_D1--D3, seed 1, 1800~s). B2 dùng 24 lần chạy (8 thử nghiệm) trên mỗi mạng. P1 và P2 với \texttt{qwen3.5:9b} chưa có kết quả chính thức nên không có mặt. Bảng~\ref{table:4.nets_results} đặt các số này cạnh mạng F2 và mạng 29/10.

\begin{table}[H]
\centering
\small
\setlength{\tabcolsep}{3pt}
\caption{B0, B1 và đặc tả B2 tốt nhất theo quy tắc dự phòng trên từng mạng (PROXY\_D1--D3, 1800~s). Mọi dòng đều bị loại bởi tiêu chí đủ điều kiện. ``Teleport kẹt/h'' gồm teleport do kẹt và do kẹt khi nhường đường.} \label{table:4.nets_results}
\begin{tabular}{|p{2.2cm}|p{3.5cm}|p{1.1cm}|p{1.1cm}|p{1.6cm}|p{1.6cm}|p{1.6cm}|}
\hline
\textbf{Mạng} & \textbf{Cấu hình} & \textbf{Số lần chạy} & \textbf{$S$} & \textbf{Va chạm/h} & \textbf{Trên làn/h} & \textbf{Teleport kẹt/h} \\ \hline
F2 & B0 & 3 & 1,181 & 112,7 & 2,0 & 10,7 \\ \hline
F2 & B1 & 3 & 1,216 & 107,3 & 0,0 & 45,3 \\ \hline
F2 & B2 (\texttt{B2\_t010}) & 3 & 1,374 & 166,0 & 46,0 & 12,0 \\ \hline
28/10 đợt 1 & B0 & 3 & 1,124 & 34,7 & 2,7 & 0,0 \\ \hline
28/10 đợt 1 & B1 & 3 & 1,198 & 41,3 & 0,0 & 24,0 \\ \hline
28/10 đợt 1 & B2 (\texttt{B2\_n28\_t000}) & 3 & 1,160 & 104,7 & 51,3 & 0,0 \\ \hline
28/10 đợt 2 & B0 & 3 & 1,212 & 38,7 & 1,3 & 7,3 \\ \hline
28/10 đợt 2 & B1 & 3 & 1,254 & 21,3 & 0,0 & 3,3 \\ \hline
28/10 đợt 2 & B2 (\texttt{B2\_n28b\_t007}) & 3 & 1,282 & 32,7 & 9,3 & 0,0 \\ \hline
29/10 & B0 & 15 & 1,225 & 50,3 & 18,5 & 1,6 \\ \hline
29/10 & B1 & 15 & 1,279 & 53,1 & 3,9 & 51,7 \\ \hline
29/10 & B2 & \multicolumn{5}{l|}{không chạy (chiến dịch chuyển sang mạng 04/11)} \\ \hline
\end{tabular}
\end{table}

\textbf{Quan sát.} Không đặc tả nào đủ điều kiện trên hai mạng 28/10 (8/8 đặc tả B2 của mỗi mạng bị loại; B0 và B1 cũng bị loại). Đặc tả B2 chọn theo quy tắc dự phòng hết teleport kẹt ở cả hai mạng, nhưng va chạm trên làn là 51,3 và 9,3/h, so với 2,7 và 1,3/h của B0. Theo $S$, 1/8 đặc tả B2 trên đợt 1 và 2/8 đặc tả trên đợt 2 có $S$ nhỏ hơn B0; giá trị nhỏ nhất là 1,091 và 1,103. Trên mạng F2 không đặc tả B2 nào có $S$ nhỏ hơn B0. Trên đợt 2, B1 có va chạm thấp nhất bảng nhưng trung bình 394 xe chưa chèn được mỗi lần chạy, và cả ba lần chạy không qua G0 (GEH 0,767--0,800).

\textbf{Diễn giải.} Mạng, tệp tuyến và cấu hình đổi cùng lúc giữa các dòng, nên chênh lệch giữa hai mạng không quy được cho một thay đổi riêng. Mỗi ngày chỉ có một seed, nên các số trên 3 lần chạy chỉ là chỉ báo. Va chạm thấp đi cùng nhiều xe chưa chèn được không phải là dấu hiệu cấu hình tốt hơn.

\subsection{Kết quả tham chiếu B0, B1 trên mạng 29/10}
\label{subsection:ref_n29}

\textbf{Thiết lập.} Ngày 29/10/2026, làn rẽ phải được giới hạn ở 30~m cuối và thêm các kết nối chỉ cho xe vi phạm: rẽ trái từ lane 0 ở hai nhánh Lê Văn Lương, quay đầu từ lane 0 và lane 1 ở cả bốn nhánh, theo đèn rẽ trái (Bảng~\ref{table:3.0}). B0 và B1 chạy trên mạng này với PROXY\_D1--D3, seed 1--5, mỗi lần 1800~s, bật kiểm tra va chạm trong nút, tắt SSM: mỗi cấu hình 15 lần chạy (Bảng~\ref{table:4.ref_n29}). B1 tại seed 1 có 1 xe rẽ trái từ lane 0 mỗi ngày ở PROXY\_D1 và D2, và 6--14 xe quay đầu từ lane 0, 1 mỗi ngày; B0 không có xe nào đi các kết nối mới.

\begin{table}[H]
\centering
\small
\caption{B0 và B1 trên mạng 29/10 (PROXY\_D1--D3 $\times$ seed 1--5, 15 lần chạy mỗi cấu hình, 1800~s). Các tỷ lệ mỗi giờ ghi trung bình $\pm$ độ lệch chuẩn giữa các lần chạy. Giá trị trong ngoặc ở ba dòng M9 là trung bình của dữ liệu đại diện.} \label{table:4.ref_n29}
\begin{tabular}{|p{5.7cm}|p{4.0cm}|p{4.0cm}|}
\hline
\textbf{Đại lượng} & \textbf{B0} & \textbf{B1} \\ \hline
Va chạm/h & 50,27 $\pm$ 9,13 & 53,07 $\pm$ 9,00 \\ \hline
Va chạm trên làn/h & 18,53 $\pm$ 6,61 & 3,87 $\pm$ 3,34 \\ \hline
Va chạm trong nút/h & 31,73 $\pm$ 9,00 & 49,20 $\pm$ 8,44 \\ \hline
Teleport kẹt/h & 1,60 $\pm$ 4,22 & 51,73 $\pm$ 48,08 \\ \hline
Xe chưa chèn cuối lần chạy (trung bình; nhỏ nhất--lớn nhất) & 1; 0--15 & 564; 352--1089 \\ \hline
Tỷ lệ ô GEH $<5$ (nhỏ nhất--lớn nhất) & 0,933--1,000 & 0,600--0,800 \\ \hline
Số lần chạy qua cổng G0 & 15/15 & 0/15 \\ \hline
Số lần chạy đủ điều kiện & 0/15 & 0/15 \\ \hline
$S$ ($S_{\text{sourced}} + S_{\text{assumed}}$) & 1,225 (0,475 + 0,750) & 1,279 (0,492 + 0,787) \\ \hline
M9a, xe/chu kỳ (đại diện 0,57) & 8,18 & 4,64 \\ \hline
M9b, tỷ lệ xe rẽ phải bị chặn (đại diện 0,065) & 0,61 & 0,80 \\ \hline
M9b, độ trễ xe rẽ phải bị chặn, s (đại diện 18,8) & 56,4 & 246,9 \\ \hline
\end{tabular}
\end{table}

\textbf{Quan sát.} B0 bị loại ở cả 15 lần chạy vì va chạm trên làn vượt 4/h (10--34/h mỗi lần chạy). Ngoài ra 2 lần có teleport không do va chạm, 1 lần còn xe chờ chèn. B1 bị loại ở cả 15 lần chạy: 15 lần còn xe chờ chèn và không qua G0, 14 lần có teleport không do va chạm, 10 lần có hơn 1 va chạm liên quan can thiệp, 5 lần vượt ngưỡng va chạm trên làn. Teleport kẹt của B1 theo seed 1--5 là 6, 3, 0, 63, 4 (PROXY\_D1); 48, 7, 1, 33, 21 (PROXY\_D2); 63, 7, 51, 28, 53 (PROXY\_D3). 8/15 lần chạy có từ 20 trở lên: nút khóa ở một số lần chạy, không phải kẹt đều ở mọi lần chạy. M9a và M9b chạm trần $D_{\text{empty}}$ ở cả B0 và B1.

\textbf{Diễn giải.} Hành vi 6 xuất hiện ở B0 mà không cần luật TraCI, khác với kết luận của thí nghiệm 3A trên mạng cấm đi thẳng từ lane 0 (Mục~\ref{subsection:phase3}). Điểm M9 chưa đọc được như sai lệch hành vi, vì hai mục tiêu của nó được suy từ giả định có cọc (Mục~\ref{section:3.6}). B1 có ít va chạm trên làn hơn B0 nhưng để lại trung bình 564 xe chưa chèn, nên hai cấu hình không so được chỉ qua số va chạm. Ngưỡng 4 va chạm trên làn/h đóng băng trên mạng F2 không còn phản ánh B0 trên mạng 29/10, và chưa được tính lại.

\subsection{Chẩn đoán khóa nút của B1 trên mạng 29/10}
\label{subsection:gridlock}

\textbf{Câu hỏi.} Teleport kẹt của B1 đến từ các kết nối mới hay từ luật? \textbf{Thiết lập.} B1 trên PROXY\_D2, 1800~s, đổi một yếu tố mỗi lần (Bảng~\ref{table:4.gridlock}).

\begin{table}[H]
\centering
\small
\caption{Biến thể chẩn đoán của B1 trên PROXY\_D2. Số trong ô ghi theo seed 1, 2, 3; biến thể chỉ chạy seed 1 có một số.} \label{table:4.gridlock}
\begin{tabular}{|p{6.4cm}|p{2.2cm}|p{2.2cm}|p{2.4cm}|}
\hline
\textbf{Biến thể} & \textbf{Teleport kẹt} & \textbf{Va chạm} & \textbf{Xe chưa chèn cuối} \\ \hline
Mạng 29/10 & 48, 7, 1 & 21, 30, 28 & 459, 451, 368 \\ \hline
Bỏ mọi kết nối mới từ lane 0, 1 & 58 & 56 & 465 \\ \hline
Chỉ bỏ rẽ trái từ lane 0 & 41 & 45 & 414 \\ \hline
Chỉ bỏ quay đầu từ lane 0, 1 & 13 & 14 & 279 \\ \hline
Không tách 30~m, không kết nối mới & 3, 16, 3 & 13, 23, 12 & 290, 370, 313 \\ \hline
Mạng 29/10, tắt luật vượt đỏ ($p_0 = p_{\max} = 0$) & 0, 0, 0 & 31, 22, 29 & 9, 98, 50 \\ \hline
\end{tabular}
\end{table}

\textbf{Quan sát.} Ở seed 1, nút khóa khoảng $t = 1290$~s: khoảng 38 xe đứng yên trong nút, riêng làn nội bộ quay đầu của lane 2 Hoàng Đạo Thúy có 7 xe vi phạm rẽ trái, và hàng chờ dồn ngược ra cả ba làn của hai nhánh Lê Văn Lương. Tắt luật vượt đỏ làm teleport kẹt về 0 ở cả 3 seed và số xe chưa chèn giảm còn 9--98, trong khi tổng va chạm giữ cùng mức. Bỏ các kết nối mới không làm hết khóa nút (58 so với 48 ở seed 1).

\textbf{Diễn giải.} Trong 3 seed này, luật vượt đỏ là điều kiện cần cho các lần khóa nút. Ba seed chưa đủ để tách phần đóng góp của mạng 29/10: mạng không tách cạnh vẫn có 16 teleport kẹt ở seed 2. Có ba cách xử lý: coi khóa nút là thuộc tính của mô hình, thêm điều kiện để xe vượt đỏ không vào nút khi lối ra kín, hoặc để P1, P2 tự chỉnh $p_0$. Ngày 04/11/2026 cách thứ hai được chọn (Mục~\ref{section:3.3}); kết quả ở Mục~\ref{subsection:ref_n1104}.

\subsection{Va chạm trên làn của B0 sau khi tách cạnh 30~m}
\label{subsection:split_collisions}

\textbf{Quan sát.} B0 trên PROXY\_D1, seed 1--5, chạy trực tiếp bằng \texttt{sumo} với cùng tùy chọn: mạng 28/10 đợt 2 có 8 va chạm trên làn và 88 va chạm trong nút; mạng 29/10 có 50 và 90. B0 không có xe đi các kết nối mới, và kết nối đi thẳng giữa hai mạng giống hệt nhau, nên khác biệt có tác dụng là điểm tách cạnh. Trong 139 va chạm trên làn của B0 ở 15 lần chạy, 105 nằm trên lane 0 và lane 1 của bốn cạnh ra, 64 là va chạm ngang.

\textbf{Diễn giải.} Cơ chế chưa xác định. Giả thuyết ban đầu là việc đổi làn bị dồn vào cạnh vạch dừng chỉ dài 30~m. Mạng 04/11 vẫn tách cạnh 30~m nhưng bỏ các kết nối từ lane 0, 1, và va chạm trên làn của B0 giảm từ 18,53/h xuống 1,87/h trên 15 lần chạy (Bảng~\ref{table:4.ref_n1104}). Giữa hai mạng chỉ có thay đổi này. Vì vậy các kết nối chỉ cho xe vi phạm có ảnh hưởng tới va chạm trên làn của B0 dù không xe nào của B0 đi chúng, và kết luận ở phần quan sát rằng khác biệt đến từ điểm tách cạnh là không đúng.

\subsection{Xe quay đầu chặn xe sau}
\label{subsection:uturn}

\textbf{Bối cảnh.} Khảo sát hiện trường ngày 30/10/2026 ghi nhận xe máy quay đầu dừng chờ làm xe sau cùng làn bị chặn, rõ nhất ở nhánh Hoàng Đạo Thúy phía bắc. Ngoài thực tế phần lớn xe máy lách sang bên hoặc sang làn kề, ô tô ít khi đổi làn để vượt.

\textbf{Đo trên mạng 29/10.} B1, PROXY\_D1, seed 1, lấy mẫu mỗi giây. Xe bị chặn là xe không quay đầu, cùng làn, cách xe quay đầu đang đứng dưới 12~m, đứng yên trong khi liên kết của chính nó xanh (Bảng~\ref{table:4.uturn_block}). B0 cũng có hiện tượng này (1207 xe$\cdot$s và 79 xe máy ở nhánh Hoàng Đạo Thúy phía bắc), nên nó không do các kết nối mới tạo ra.

\begin{table}[H]
\centering
\small
\caption{Thời gian xe bị chặn sau xe quay đầu đang đứng (B1, PROXY\_D1, seed 1, xe$\cdot$s).} \label{table:4.uturn_block}
\begin{tabular}{|p{3.2cm}|p{2.6cm}|p{2.8cm}|p{1.6cm}|p{2.2cm}|}
\hline
\textbf{Nhánh} & \textbf{Xe máy đi thẳng} & \textbf{Xe vi phạm rẽ trái} & \textbf{Ô tô} & \textbf{Số xe máy khác nhau} \\ \hline
HDT\_N & 1409 & 3053 & 0 & 79 \\ \hline
HDT\_S & 285 & 833 & 22 & 101 \\ \hline
LVL\_E & 658 & 615 & 209 & 47 \\ \hline
LVL\_W & 1720 & 111 & 231 & 83 \\ \hline
\end{tabular}
\end{table}

\textbf{Luật và kết quả.} Đặc tả \texttt{B1\_uturn\_bypass} là B1 giữ nguyên, thêm luật \texttt{uturn\_bypass} (Bảng~\ref{table:3.2}) với $p = \tamthoi{0{,}8}$ cho xe máy, $p = \tamthoi{0{,}1}$ cho ô tô, $t_{\text{wait}} = \tamthoi{2}$~s; các giá trị này là giả định từ mô tả định tính. Trên kịch bản nhỏ luật chạy đúng: một ô tô đứng sau xe quay đầu đổi làn lúc 58,4~s và qua vạch lúc 63,2~s, so với 78,1~s khi không có luật. Bảng~\ref{table:4.uturn_rule} cho kết quả trên mạng đại diện qua ba phiên bản. Lần B1 cùng ngày, cùng seed có 6 teleport kẹt.

\begin{table}[H]
\centering
\small
\caption{Luật \texttt{uturn\_bypass} trên PROXY\_D1, seed 1, 1800~s. v1: chỉ lệnh lách hoặc đổi làn, lề 0,3~m. v2: thêm lane change mode 01. v3: thêm \texttt{lcPushy} bằng 1, lề 0,1~m.} \label{table:4.uturn_rule}
\begin{tabular}{|p{1.4cm}|p{1.9cm}|p{1.9cm}|p{2.6cm}|p{2.6cm}|p{1.9cm}|}
\hline
\textbf{Phiên bản} & \textbf{Số lần rút trúng} & \textbf{Vượt được} & \textbf{Không có khe hoặc làn} & \textbf{Không vượt trong 6~s} & \textbf{Teleport kẹt} \\ \hline
v1 & 144 & 8 & 88 & 48 & 7 \\ \hline
v2 & 150 & 16 & 77 & 57 & 20 \\ \hline
v3 & 150 & 13 & 42 & 95 & 28 \\ \hline
\end{tabular}
\end{table}

\textbf{Diễn giải.} Tỷ lệ vượt được là 13/150 ở phiên bản cuối, nên luật có tác dụng nhỏ trên mạng đại diện. Trong hàng xe máy dày, SUMO giữ hình học không chồng lấn: ba hàng xe đứng song song trên làn 3,5~m, khoảng cách dọc 1,9~m (\texttt{minGap} của B1) ngắn hơn chiều dài xe 2,2~m, nên xe sau không có khe để lách. Teleport kẹt tăng từ 6 lên 28 trên cùng một seed chưa quy được cho luật, vì ở Bảng~\ref{table:4.ref_n29} cùng một đặc tả B1 có 0--63 teleport kẹt tùy seed. Muốn có hành vi lách qua cần \texttt{minGap} nhỏ hơn cho xe máy đứng chờ, hoặc cho xe vi phạm rẽ trái đứng lệch sang một bên bằng \texttt{latAlignment}; cả hai đều đổi cấu trúc B1 nên chưa làm. Ngày 04/11/2026 luật được đưa vào cấu trúc B1 với các giá trị trên và vào prompt từ \texttt{v003} (Mục~\ref{subsection:ref_n1104}).

\subsection{Thay đổi ngày 04/11 và kết quả tham chiếu trên mạng 04/11}
\label{subsection:ref_n1104}

\textbf{Thay đổi.} Ngày 04/11/2026 có ba thay đổi, chọn sau các chẩn đoán ở Mục~\ref{subsection:gridlock} đến~\ref{subsection:uturn}: (i) mạng bỏ các kết nối chỉ cho xe vi phạm từ lane 0, 1 (rẽ trái từ lane 0, quay đầu từ lane 0 và lane 1), nên rẽ trái và quay đầu chỉ còn từ lane 2, và mạng này gọi là mạng 04/11; (ii) luật \texttt{swarm\_red\_running} chỉ cho xe vào nút khi nút thông (Mục~\ref{section:3.3}); và (iii) \texttt{uturn\_bypass} vào cấu trúc cố định của B1 và vào prompt \texttt{v003}.

\textbf{Thiết lập.} B0 và B1 trên mạng 04/11, PROXY\_D1--D3 $\times$ seed 1--5, mỗi lần 1800~s, bật kiểm tra va chạm trong nút; mỗi cấu hình 15 lần chạy (chiến dịch \texttt{REF\_n1104s5}). Bảng~\ref{table:4.ref_n1104} đặt kết quả cạnh mạng 29/10 (Bảng~\ref{table:4.ref_n29}).

\begin{table}[H]
\centering
\small
\caption{B0 và B1 trên mạng 29/10 và mạng 04/11 (PROXY\_D1--D3 $\times$ seed 1--5, 15 lần chạy mỗi ô, 1800~s; trung bình trên 15 lần chạy).} \label{table:4.ref_n1104}
\begin{tabular}{|p{4.4cm}|p{2.1cm}|p{2.1cm}|p{2.1cm}|p{2.1cm}|}
\hline
\textbf{Đại lượng} & \textbf{B0 29/10} & \textbf{B0 04/11} & \textbf{B1 29/10} & \textbf{B1 04/11} \\ \hline
Va chạm/h & 50,27 & 36,13 & 53,07 & 33,33 \\ \hline
Va chạm trên làn/h & 18,53 & 1,87 & 3,87 & 0,27 \\ \hline
Teleport kẹt/h & 1,60 & 0,00 & 51,73 & 4,67 \\ \hline
$S$ & 1,2247 & 1,2174 & 1,2791 & 1,2494 \\ \hline
Số lần chạy đủ điều kiện & 0/15 & 15/15 & 0/15 & 0/15 \\ \hline
\end{tabular}
\end{table}

\textbf{Quan sát.} B0 đủ điều kiện ở 15/15 lần chạy. Đây là lần đầu B0 qua mọi điều kiện trên mạng đại diện. B0 không có luật TraCI, nên chỉ thay đổi mạng tác động lên nó. B1 vẫn bị loại ở cả 15 lần chạy: 14 lần còn xe chờ chèn lúc kết thúc (22--457 xe); 14 lần có hơn 1 va chạm liên quan can thiệp (2--6); 9 lần có teleport không do va chạm (1--10); 5 lần không qua G0, cả 5 đều ở PROXY\_D1. Không lần nào vượt ngưỡng va chạm trên làn. Số xe chờ chèn lớn nhất ở PROXY\_D1 (241--457 xe). Ở mạng 29/10 mỗi lần nút khóa sinh 20--40 teleport kẹt; ở mạng 04/11 một lần chạy B1 có nhiều nhất 10 teleport không do va chạm. Nhật ký can thiệp của B1, cộng 15 lần chạy: \texttt{swarm\_red\_running} rút trúng 3380 lượt vượt đỏ, thực hiện 1946 lượt; 2775 lượt tới vạch khi nút chưa thông nên phải đợi, và 1016 lượt bị hủy vì đèn hết đỏ trước khi nút thông. \texttt{uturn\_bypass} có 1927 lượt thử, 116 lượt lách qua được, 494 lượt không có khe ngang và không có làn kề, 1317 lượt không vượt được xe quay đầu. \texttt{early\_start} và \texttt{queue\_space\_seeking} không can thiệp lần nào vì B1 đặt mọi xác suất bằng 0.

\textbf{Diễn giải.} Khoảng 30\% số lượt vượt đỏ đã rút không thành hành vi trong mô phỏng, nên khi hiệu chỉnh $p_0$, $k$, $p_{\max}$, tỷ lệ vượt đỏ đo được (M3) sẽ thấp hơn tỷ lệ rút. B1 có hai thay đổi cùng lúc (điều kiện nút thông và \texttt{uturn\_bypass}), nên bảng trên chưa tách được tác dụng riêng của từng thay đổi. Mục~\ref{subsection:split_collisions} ghi khác biệt về va chạm trên làn của B0 giữa hai mạng.

\subsection{Luật cắt về cho xe đi sai làn}
\label{subsection:cutback}

\textbf{Bối cảnh.} Bỏ các kết nối từ lane 0, 1 làm mất hành vi xe vi phạm rẽ trái từ lane 0 và quay đầu từ lane 0, 1. Hành vi này được làm lại bằng lựa chọn \texttt{wrong\_lane} kiểu cắt về của luật \texttt{queue\_space\_seeking} (Mục~\ref{section:3.3}); mạng không đổi và luật không thêm tham số.

\textbf{Thí nghiệm trước khi viết luật.} PROXY\_D1, B1 không có luật TraCI, 1200~s, 12 xe được ép đi làn sai. Nếu chỉ cho xe cắt về khi đã tới vạch và đèn hướng của nó xanh, thì đầu lane 2 đã kín: xe phải đứng ở cuối làn sai 1,5--8 phút, và một xe bị teleport. Tăng \texttt{lcMaxSpeedLatStanding} lên 3~m/s không đổi kết quả. Vì vậy luật cho xe nhập lại ngay khi vào đoạn 30~m cuối.

\textbf{Chạy thử.} Đặc tả \texttt{B1\_cut\_back} là B1, thêm \texttt{queue\_space\_seeking} cho \texttt{moto\_seeker} và \texttt{car\_seeker} với $p_{\text{wrong\_lane}} = 0{,}5$; PROXY\_D1--D3 $\times$ seed 1.
Có 49 lượt rút \texttt{wrong\_lane}. 15 lượt vào được làn sai rồi cắt về: 10 lượt quay đầu từ lane 1 (7 xe máy, 3 ô tô) và 5 lượt từ lane 0. 31 lượt không đổi được làn trong 5~s cho mỗi làn phải vượt; 3 lượt rút khi xe đã ở trong đoạn 30~m. 14 trong 15 xe cắt về đi hết hành trình. Một xe bị teleport vì chờ nhường quá lâu ở lane 2, 13 phút sau khi cắt về, trong lúc tắc chung. Một ô tô có va chạm ở lane 2, 40~s sau khi cắt về. Cả 3 lần chạy bị loại: 37,3 va chạm/h, 2,0 va chạm trên làn/h, 9,3 teleport kẹt/h, $S = 1{,}2636$. B1 cùng seed (Bảng~\ref{table:4.ref_n1104}) có 42,7 va chạm/h, 0 va chạm trên làn/h, 5,3 teleport kẹt/h. Ba lần chạy không đủ để nói luật làm tăng hay giảm các số này.

\textbf{Đo M8.} Xe cắt về vào nút từ lane 2, nên M8 đọc tại vạch dừng không ghi nhận nó: M8 của \texttt{B1\_cut\_back} bằng 0 như B1, trong khi trung bình của dữ liệu đại diện là 0,249 với xe máy và 0,068 với ô tô. Ngày 08/11/2026, M8 được chuyển sang đọc làn ở đầu đoạn 30~m (Mục~\ref{section:3.7}). M8 theo định nghĩa mới chưa được tính lại cho ba lần chạy này.

\subsection{Rút gọn prompt cho \texttt{qwen3.5:9b}}
\label{subsection:prompt_v004}

\textbf{Vấn đề.} Ngày 04/11, với prompt \texttt{v003}, \texttt{qwen3.5:9b} cho 0/3 chuỗi thử có đặc tả hợp lệ sau 3 lần thử; ngày 03/11 với \texttt{v002} là 1/3. Có hai nguyên nhân chính. Thứ nhất, nguyên mẫu không ghi \texttt{distribution} rơi vào phân phối mặc định của lớp xe, và hai nguyên mẫu \texttt{custom1} cùng rơi vào \texttt{moto\_violator}, làm sai tổng tỷ trọng. Thứ hai, prompt quá dài: dựng lại từ chiến dịch P2 trên mạng 28/10, prompt vòng 3 khoảng 46.300 token, vượt cửa sổ 32.768 token (Bảng~\ref{table:4.prompt_size}). Prompt \texttt{v004} sửa hai điểm này (Bảng~\ref{table:3.prompts}).

\begin{table}[H]
\centering
\small
\caption{Kích thước prompt P2 dựng lại từ kết quả chiến dịch P2 trên mạng 28/10 đợt 2. Số token ước theo tỷ lệ ký tự trên token đo ngày 03/11.} \label{table:4.prompt_size}
\begin{tabular}{|p{2.0cm}|p{5.0cm}|p{5.0cm}|}
\hline
\textbf{Vòng} & \textbf{\texttt{v002}} & \textbf{\texttt{v004}} \\ \hline
Vòng 2 & 68.745 ký tự, khoảng 25.200 token & 45.029 ký tự, khoảng 16.500 token \\ \hline
Vòng 3 & 126.372 ký tự, khoảng 46.300 token & 58.029 ký tự, khoảng 21.300 token \\ \hline
\end{tabular}
\end{table}

\textbf{Thử trên máy chính.} Request là prompt vòng 2 của P2 dựng bằng \texttt{v004}; thinking bật, \texttt{max\_tokens} 4500; mỗi chuỗi tối đa 3 lần thử, seed 42, 43, 44; hai lượt, lượt sau có đủ các thay đổi của \texttt{v004}. Cả 6 chuỗi đều cho đặc tả hợp lệ trong 3 lần thử. Prompt dài 15.919--19.678 token, câu trả lời 1.526--4.500 token, tốc độ sinh 12,5--13,5 token/s, mỗi lần thử 2--6 phút. Lỗi đặt mã nguyên mẫu trùng tên phân phối \texttt{moto\_violator} xuất hiện ở 4/6 chuỗi dù prompt đã cấm, và lần thử sau đều sửa được. Hai lần thử bị cắt ở 4500 token. Sáu chuỗi không đủ để so hai lượt với nhau.

\subsection{Chiến dịch B2, P1, P2 trên mạng 04/11}
\label{subsection:campaigns_n1108}

\textbf{Thiết lập.} PROXY\_D1--D3 $\times$ seed 1, ngân sách 36 lần chạy mỗi chiến dịch, mỗi lần 1800~s. Mức trần va chạm trên làn khi chọn đặc tả tốt nhất là \tamthoi{1,8667} $+ 1$ va chạm/h, với \tamthoi{1,8667} là B0 của \texttt{REF\_n1104s5}. B2 là Optuna với 3 tiến trình. P1 (\texttt{params\_only}, nền B1) và P2 (\texttt{full}) dùng \texttt{qwen3.5:9b}, prompt \texttt{v004}, 4 mẫu mỗi vòng, tối đa 3 lần thử mỗi mẫu, \texttt{max\_tokens} 7000, thinking tắt. Ba chiến dịch chạy ngày 08/11/2026; B2 chạy song song với P1, P2 chạy sau P1.

Thinking bị tắt vì hai lần khởi động P1 đầu tiên đều hỏng ngay lần thử đầu của vòng 1: câu trả lời bị cắt ở 4500 rồi ở 7000 token, phần thinking dài 18.554 và 28.585 ký tự, chưa tới JSON. Hai thư mục chiến dịch đó bị xóa trước khi có lần chạy SUMO nào được tính; nhật ký LLM của chúng được giữ. Với thinking tắt, chạy thử trên prompt vòng 1 cho 3/3 chuỗi P1 và 2/3 chuỗi P2 hợp lệ.

\begin{table}[H]
\centering
\small
\caption{Ba chiến dịch trên mạng 04/11 (PROXY\_D1--D3 $\times$ seed 1, 36 lần chạy mỗi chiến dịch). Giờ Hà Nội ngày 08/11/2026.} \label{table:4.campaigns_n1108}
\begin{tabular}{|p{3.6cm}|p{3.2cm}|p{3.2cm}|p{3.2cm}|}
\hline
 & \textbf{B2} & \textbf{P1} & \textbf{P2} \\ \hline
Thời gian & 3~h~25 (06:03--09:28) & khoảng 2~h~25 (tới 09:20) & 1~h~56 (09:20--11:16) \\ \hline
Thử nghiệm hoặc vòng & 12 thử nghiệm & 4 vòng (3 + 1) & 5 vòng (3 + 2) \\ \hline
Mẫu hợp lệ & 12/12 & 12/16 & 14/20 (2 mẫu hợp lệ nhưng hết ngân sách) \\ \hline
Lần trả lời hợp lệ & không dùng LLM & 12/27 & 14/36 \\ \hline
Đặc tả được chạy & 12 & 12 & 12 \\ \hline
Đặc tả đủ điều kiện & 0 & 0 & 0 \\ \hline
Lần chạy đủ điều kiện PROXY\_D1 / D2 / D3 & 0/12, 0/12, 0/12 & 0/12, 0/12, 0/12 & 0/12, 6/12, 5/12 \\ \hline
$S$ nhỏ nhất & 1,1438 (\texttt{t003}) & 1,2386 (\texttt{it01\_s01}) & 1,1680 (\texttt{it04\_s03}) \\ \hline
Đặc tả chọn theo quy tắc dự phòng & \texttt{t002}, $S = 1{,}2088$ & \texttt{it01\_s02}, $S = 1{,}2984$ & \texttt{it03\_s02}, $S = 1{,}1744$ \\ \hline
\end{tabular}
\end{table}

\begin{sloppypar}
\textbf{Quan sát.} Không chiến dịch nào có đặc tả đủ điều kiện, nên đặc tả tốt nhất được chọn theo quy tắc dự phòng: ít teleport kẹt nhất, rồi theo $S$, không xét mức trần va chạm. B2 và P1 còn xe chờ chèn lúc kết thúc ở cả ba ngày. P1 có ở 36/36 lần chạy (PROXY\_D1 301--623 xe, PROXY\_D2 23--179, PROXY\_D3 133--297). B2 có ở 21/36 lần chạy, trong đó cả 12 lần chạy PROXY\_D1 (20--338 xe). Gần như mọi đặc tả có lần chạy với hơn 1 va chạm liên quan can thiệp. 9/12 thử nghiệm B2 còn vượt ngưỡng va chạm trên làn ở ít nhất một lần chạy; P1 không có lần chạy nào bị loại vì lý do này. P2 chỉ còn xe chờ chèn ở PROXY\_D1 (4--163 xe, cả 12 đặc tả). Ở PROXY\_D2 và PROXY\_D3, 11/24 lần chạy đủ điều kiện; các lần còn lại bị loại vì va chạm trên làn vượt ngưỡng hoặc vì hơn 1 va chạm liên quan can thiệp. Từ vòng 2, cả 12 đặc tả P2 được chạy có cùng cấu trúc với \texttt{it02\_s01}: 5 nguyên mẫu (\texttt{car\_compliant}, \texttt{car\_early\_starter}, \texttt{moto\_compliant}, \texttt{moto\_slow\_center} và \texttt{moto\_violator\_left} lớp \texttt{custom1}) và 3 luật (\texttt{early\_start} cho ô tô, \texttt{swarm\_red\_running} cho xe máy thường, \texttt{illegal\_left\_box}). Các vòng sau chỉ đổi tỷ trọng và tham số. Cấu trúc này không có \texttt{queue\_space\_seeking}, \texttt{uturn\_bypass} (chỉ \texttt{it01\_s02} có), \texttt{slow\_lane\_center} (dù có nguyên mẫu tên \texttt{moto\_slow\_center}), và không có xe đi sai làn.
\end{sloppypar}

\textbf{Diễn giải.} P1 giữ cấu trúc của B1, mà B1 đã có 14/15 lần chạy còn xe chờ chèn (Mục~\ref{subsection:ref_n1104}), nên P1 không thể cho đặc tả đủ điều kiện dù LLM chọn tham số thế nào. $S$ thấp hơn của P2 đến từ đặc tả bỏ bớt một số hành vi của kịch bản. Ba chiến dịch chưa tách được phần giảm xe chờ chèn của P2 là do tham số vType hay do bỏ các luật có thể chặn làn (cắt về, lách qua xe quay đầu). Sau kết quả này, bộ kiểm tra hợp lệ bắt đặc tả P2 dùng đủ sáu luật từ prompt \texttt{v005} (Mục~\ref{section:3.4}). Các chiến dịch chỉ dùng seed 1, nên $S$ ở Bảng~\ref{table:4.campaigns_n1108} không so trực tiếp được với $S$ của B0 và B1 trên 5 seed ở Bảng~\ref{table:4.ref_n1104}.

\subsection{Hành vi của \texttt{qwen3.5:9b} trong chiến dịch}
\label{subsection:qwen_campaign}

Prompt dài 9.471--24.441 token ở P1 và 7.063--19.895 token ở P2. Câu trả lời dài 1.257--2.256 token, và cả 63 câu trả lời kết thúc bình thường, không câu nào bị cắt. Mỗi lần thử mất khoảng 1,5--3,5 phút.

\textbf{P1.} 14/15 lần thử hỏng là lỗi tổng tỷ trọng, cả 14 lần ở phân phối \texttt{moto\_violator}. Phân phối này chỉ có một nguyên mẫu (\texttt{moto\_violator\_saval}), nên tỷ trọng của nó phải bằng 1, nhưng LLM đổi thành 0,9, 0,85, 0,5, 0,1 hoặc 0,05. Ở 2 lần thử của \texttt{it03\_s03}, LLM chuyển 0,1 và 0,05 tỷ trọng từ \texttt{moto} sang \texttt{moto\_violator}, nên cả hai phân phối đều sai. 4 mẫu hỏng cả 3 lần thử vì lỗi này, dù tin nhắn thử lại đã nêu rõ tổng tỷ trọng. P1 còn nộp lại một đặc tả: \texttt{it01\_s02}, \texttt{it02\_s03}, \texttt{it03\_s01} và \texttt{it04\_s02} giống nhau (\texttt{early\_start} $p = 0{,}13$, $p_{\text{sidewalk}} = 0{,}75$, $p_{\text{wrong\_lane}} = 0{,}25$, phần còn lại như B1). Đây là đặc tả được chọn sau vòng 1, nên phản hồi giới thiệu nó là tốt nhất. Ba lần nộp lại tốn 9 lần chạy và cho cùng $S = 1{,}2984$.

\textbf{P2.} Ở vòng 1 và 2, 6/8 mẫu không hợp lệ sau 3 lần thử. Bảng~\ref{table:4.qwen_errors} đếm lỗi trên mọi lần thử hỏng. Từ vòng 3, 11/12 mẫu hợp lệ ngay lần thử đầu. Từ vòng 2, lịch sử gửi cho LLM có JSON của một đặc tả hợp lệ (\texttt{it02\_s01}), và các đặc tả vòng 3--5 giữ đúng cấu trúc đó. Mọi lỗi trong bảng là lỗi cấu trúc, và chúng không xuất hiện lại khi LLM giữ nguyên một cấu trúc đã qua bộ kiểm tra.

\begin{table}[H]
\centering
\small
\caption{Lỗi trên các lần thử hỏng của P2 (một lần thử có thể có nhiều lỗi).} \label{table:4.qwen_errors}
\begin{tabular}{|p{11.0cm}|p{1.8cm}|}
\hline
\textbf{Lỗi} & \textbf{Số lần} \\ \hline
Mã nguyên mẫu trùng tên phân phối (\texttt{car}, \texttt{moto}, \texttt{moto\_violator}, \texttt{car\_violator}) & 19 \\ \hline
Tổng tỷ trọng của một phân phối khác 1 & 15 \\ \hline
Phân phối mà nhu cầu dùng không có nguyên mẫu nào & 7 \\ \hline
\texttt{queue\_space\_seeking} áp cho lớp xe không được phép & 6 \\ \hline
\texttt{queue\_space\_seeking} không áp dụng được trên mạng cho nguyên mẫu trong \texttt{moto\_violator} & 5 \\ \hline
\texttt{illegal\_left\_box} áp cho xe máy thường & 4 \\ \hline
Giá trị ngoài miền (lược đồ) & 4 \\ \hline
\texttt{early\_start} áp hai lần cho cùng nguyên mẫu & 1 \\ \hline
\end{tabular}
\end{table}

\subsection{Nguyên nhân xe chờ chèn ở PROXY\_D1}
\label{subsection:d1_backlog}

\textbf{Câu hỏi.} Ở PROXY\_D1, mọi đặc tả của ba chiến dịch và 5/5 lần chạy B1 còn xe chờ chèn, trong khi B0 với cùng nhu cầu thì không. Phần nào của B1 gây ra điều này?

\textbf{Thiết lập.} B1 trên PROXY\_D1 $\times$ seed 1--3, 1800~s. Bỏ lần lượt từng luật TraCI, rồi bỏ mọi luật và đưa từng nhóm tham số vType về giá trị của B0 (mặc định SUMO). Trong B1, \texttt{early\_start} và \texttt{queue\_space\_seeking} có mọi xác suất bằng 0 và không can thiệp lần nào, nên ở các biến thể bỏ một luật chỉ \texttt{swarm\_red\_running} và \texttt{uturn\_bypass} được tắt.

\begin{table}[H]
\centering
\small
\caption{Biến thể của B1 trên PROXY\_D1. Số xe chờ chèn ghi theo seed 1, 2, 3.} \label{table:4.d1_backlog}
\begin{tabular}{|p{6.4cm}|p{3.0cm}|p{1.8cm}|p{1.4cm}|}
\hline
\textbf{Biến thể} & \textbf{Xe chờ chèn lúc kết thúc} & \textbf{Va chạm/h} & \textbf{$S$} \\ \hline
B1 đủ luật & 339, 457, 377 & 24--26 & 1,2712 \\ \hline
Bỏ \texttt{swarm\_red\_running} & 332, 337, 221 & 14--18 & 1,2638 \\ \hline
Bỏ \texttt{uturn\_bypass} & 295, 273, 291 & 20--32 & 1,2783 \\ \hline
Bỏ mọi luật & 387, 203, 227 & 12--20 & 1,2974 \\ \hline
Bỏ mọi luật, ô tô dùng tham số B0 & 295, 206, 263 & 24--32 & 1,2632 \\ \hline
Bỏ mọi luật, xe máy bỏ \texttt{accel} (mặc định 6~m/s\textsuperscript{2}) & 10, 126, 52 & 48--58 & 1,2483 \\ \hline
\end{tabular}
\end{table}

\textbf{Quan sát.} Bỏ mọi luật TraCI vẫn còn 203--387 xe chờ chèn. Khi trả \texttt{accel} của xe máy về mặc định, số xe chờ chèn giảm xuống 10--126 nhưng chưa về 0, và va chạm tăng từ 12--20/h lên 48--58/h. Tham số ô tô (\texttt{tau} 2,5~s, \texttt{accel} 1,2~m/s\textsuperscript{2}) chỉ làm số xe chờ chèn thay đổi ít. B2 tìm \texttt{accel} của xe máy trong $[0{,}8;\,6]$~m/s\textsuperscript{2}, nhưng cả 12 thử nghiệm vẫn có 20--338 xe chờ chèn ở PROXY\_D1; thử nghiệm \texttt{t000} với \texttt{accel} 5,8 vẫn còn 256 xe, và có \texttt{tau} xe máy 2,5~s.

\textbf{Diễn giải.} Nguyên nhân chính là tham số vType của B1, rõ nhất là \texttt{accel} 1,1~m/s\textsuperscript{2} của xe máy. Giá trị này chuyển từ gia tốc cực đại $a$ của mô hình IDM trong Bảng 11 của Saval và cộng sự~\cite{saval2023mixed} sang \texttt{accel} của mô hình Krauss (Phụ lục~\ref{section:appendix_b1}), trong khi mặc định SUMO cho xe máy là 6~m/s\textsuperscript{2}. Với \texttt{accel} thấp, xe máy tăng tốc chậm và mỗi pha xanh xả được ít xe. Chưa có biến thể nào thử riêng \texttt{tau} của xe máy. Việc giữ giá trị của Saval trong B1 hay sửa B1 (dùng mặc định SUMO cho \texttt{accel} xe máy, hoặc nâng cận dưới của \texttt{accel} xe máy trong tệp miền giá trị) chưa được quyết định.

\subsection{Các thực nghiệm dự kiến}
\label{subsection:planned_experiments}

Việc đầu tiên là quyết định về \texttt{accel} của xe máy trong B1 (Mục~\ref{subsection:d1_backlog}), vì khi PROXY\_D1 còn xe chờ chèn, không đặc tả nào dựa trên cấu trúc B1 đủ điều kiện được. Sau đó, P2 được chạy lại với prompt \texttt{v005} để đặc tả có đủ sáu luật TraCI, và đặc tả chọn được của B2 và P2 được chạy lại với 5 seed như \texttt{REF\_n1104s5} trước khi so với B0 và B1. Phương án H chạy Optuna với cùng ngân sách 36 lần chạy trên cấu trúc của đặc tả tốt nhất của P2.

Đánh giá cuối chạy B0, B1, B2, P1, P2 và H trên hai ngày đánh giá PROXY\_D4--PROXY\_D5 với ít nhất 20 seed. E1 so sánh P1 với B2; E2 so sánh B2, H và P2. Mỗi cặp dùng kiểm định Mann--Whitney có hiệu chỉnh Holm và kích thước hiệu ứng Cliff's delta, kèm một lần chạy chẩn đoán 3600~s $\times$ 3 seed cho B0 và đặc tả tốt nhất trên PROXY\_D4. Nếu còn thời gian, luận văn làm thêm phân tích loại bỏ đầy đủ (tắt lần lượt từng luật TraCI hoặc từng nhóm tham số trên mọi ngày hiệu chỉnh), và kịch bản chính sách với lưu lượng rẽ trái trái phép bằng 0\%, 50\% và 100\% mức quan sát, đo độ trễ và số xung đột. Kịch bản chính sách sẽ được ghi rõ là ngoại suy.
Hai cải tiến của vòng LLM đã xác định nhưng chưa làm: dùng lại kết quả khi LLM nộp lại một đặc tả trùng với đặc tả đã chạy (cùng giá trị băm) thay vì chạy lại; và ở chế độ \texttt{params\_only}, ép tỷ trọng bằng 1 ngay trong lược đồ cho phân phối chỉ có một nguyên mẫu (ở P1, 12/15 lần thử hỏng chỉ sai ở phân phối đó). Hai mục tiêu của M9 vẫn chưa được định nghĩa lại theo mô hình làn 30~m (Mục~\ref{section:3.6}).

\section{Triển khai và tái lập}
\label{section:deploy}

\textbf{Môi trường.} Hệ thống phát triển trên máy Windows, SUMO 1.27.0 cài sẵn, biến \texttt{SUMO\_HOME} trỏ tới thư mục cài đặt, Máy chính chạy Windows 11 Pro với CPU Intel Core i5-13400F (10 nhân, 16 luồng), 31,8~GB RAM và GPU NVIDIA GeForce RTX 4060 8~GB; LLM cục bộ chạy trên máy này. Một máy Linux phụ chạy bổ sung một phần chiến dịch, với SUMO 1.27.0 cài dưới dạng gói Python \texttt{eclipse-sumo} trong môi trường ảo. Tám lần chạy B0 trên máy chính được chạy lại trên máy phụ và cho đúng từng số va chạm và teleport. Môi trường Python là môi trường ảo Python 3.12, dựng bằng cùng một script trên cả hai máy. Gói dự án cài ở chế độ có thể chỉnh sửa kèm nhóm phụ thuộc phát triển; \texttt{libsumo} và \texttt{sumo-data} 1.27.0 cài không kèm phụ thuộc; một tệp đường dẫn trong môi trường ảo đưa thư mục công cụ của SUMO vào đường dẫn tìm mô-đun, để nạp \texttt{traci} và \texttt{sumolib}. Toàn bộ kiểm thử chạy bằng \texttt{pytest} từ thư mục gốc của dự án.

\begin{sloppypar}
\textbf{Lệnh tái lập.} Các lệnh mô-đun Python theo thứ tự của Bảng~\ref{table:4.2}: (i) \texttt{thesis\_sumo.net.build} với tệp cấu hình hiện trường đại diện: dựng mạng; (ii) \texttt{thesis\_sumo.proxy.generate}: sinh năm ngày dữ liệu đại diện và tệp tuyến; (iii) \texttt{thesis\_sumo.rules.compiler}: biên dịch một đặc tả (bắt buộc truyền tệp cấu hình đại diện, vì cấu hình đo thực chưa có bề rộng làn); (iv) \texttt{thesis\_sumo.sim.runner}: chạy mô phỏng, với tham số mã đặc tả, ngày, danh sách seed, tệp cấu hình hiện trường, thời điểm kết thúc, số tiến trình song song, chế độ SSM và tùy chọn mở \texttt{sumo-gui}; và (v) \texttt{thesis\_sumo.baselines.optuna\_search} (B2, H) và \texttt{thesis\_sumo.llm.loop} (P1 với chế độ \texttt{params\_only} và cơ sở B1, P2 với chế độ \texttt{full}): chạy chiến dịch, cùng ngân sách, cùng cấu hình đại diện, cùng mức trần va chạm của B0 trên mạng đang dùng.
\end{sloppypar}
Thí nghiệm vùng chờ giữa nút và Phase 3A--3C có script riêng. Khi dùng backend tệp, lệnh \texttt{loop pending} liệt kê yêu cầu chưa có phản hồi; chạy lại cùng lệnh sau khi có phản hồi thì chiến dịch tiếp tục. Dữ liệu mỗi chiến dịch nằm trong \texttt{experiments/\allowbreak campaigns/\allowbreak <chiến dịch>/} (ngân sách, đặc tả, đặc tả đã biên dịch, lần chạy, trạng thái vòng lặp hoặc nghiên cứu Optuna); nhật ký LLM nằm trong \texttt{experiments/\allowbreak llm\_logs/\allowbreak <chiến dịch>/}.

\textbf{Đầu ra của một lần chạy.} Mỗi lần chạy ghi vào \texttt{experiments/\allowbreak runs/\allowbreak <spec\_id>\_\_<day>/\allowbreak seed\_XX/}, với \texttt{spec\_id} là mã đặc tả, \texttt{day} là ngày dữ liệu và \texttt{XX} là seed; thư mục seed bị xóa và tạo lại nếu chạy lại. Thư mục gồm: (i) \texttt{run.sumocfg} và tệp khai báo đầu ra; (ii) đầu ra SUMO: thông tin chuyến đi, quỹ đạo nén, chuyển làn, thống kê, va chạm, trạng thái đèn, detector; (iii) \texttt{sumo.log}; \texttt{vtypes.json} (các vType SUMO đã nạp); \texttt{interventions.csv} nếu có luật TraCI; (iv) \texttt{manifest.json}: mã lần chạy, mã đặc tả, ngày, seed; mã commit và trạng thái thay đổi chưa commit; phiên bản SUMO và backend; giá trị băm của đặc tả, mạng, tệp tuyến, cấu hình hiện trường; nhãn nguồn cấu hình hiện trường; cấu hình thí nghiệm và tùy chọn SUMO; thời gian chạy; số teleport (tổng và không do va chạm); số va chạm (tổng, theo loại, trên làn, trong nút, liên quan can thiệp TraCI); tóm tắt can thiệp và thời gian xử lý của bộ điều khiển; số xe chưa chèn (cuối lần chạy và lớn nhất); số cảnh báo và lỗi; cờ \texttt{valid\_strict}; và (v) thư mục các bảng CSV theo lược đồ quan sát.
Đầu ra thô không đưa vào kho mã; kho mã chỉ lưu mã, cấu hình, tệp vá, mạng đã sinh, dữ liệu đại diện và các bảng kết quả nhỏ. Một lần chạy B0 1800~s trên dữ liệu đại diện, khi chạy 3 lần song song, mất 88--463~s.

Chương này đã mô tả phần mềm hiện thực phương pháp, bộ kiểm thử 403 ca, và các thực nghiệm qua sáu phiên bản mạng. Các thí nghiệm khả thi xác nhận SUMO tự mô phỏng được xe chờ giữa nút với sức chứa 8 xe, và chỉ ra nguồn va chạm do bề rộng làn ra. Trên mạng 04/11, sau khi bỏ các kết nối vi phạm từ lane 0, 1 và thêm điều kiện nút thông cho luật vượt đỏ, SUMO mặc định (B0) lần đầu đủ điều kiện ở 15/15 lần chạy. Ba chiến dịch B2, P1 và P2 với LLM cục bộ \texttt{qwen3.5:9b} đều chạy hết ngân sách 36 lần chạy nhưng không có đặc tả nào đủ điều kiện. Nguyên nhân chính là xe chờ chèn ở PROXY\_D1, mà phân tích loại bỏ quy cho \texttt{accel} 1,1~m/s\textsuperscript{2} của xe máy trong B1, không phải cho các luật TraCI. Ngoài ra, P2 đạt $S$ thấp hơn bằng cách bỏ bớt luật, nên từ prompt \texttt{v005} bộ kiểm tra hợp lệ bắt đặc tả P2 dùng đủ sáu luật. Mọi kết quả trên tính trên dữ liệu đại diện và các giá trị tạm thời, nên chúng cho biết quy trình vận hành được, chưa cho biết mô phỏng giống giao thông thực đến đâu. Chương~\ref{chapter:SolutionAndContribution} tổng hợp các giải pháp kỹ thuật rút ra từ quá trình này.

\end{document}
